\chapter{Att kontrollera ett påstående}
\label{app:verifiera-grafik}

Det här materialet gör, som allt undervisningsmaterial, en rad
påståenden.  De flesta går att pröva vid datorn, och det är poängen med
att programmen i texten är de program som körs: att parenteserna efter
\mintinline{python}{greet} anropar funktionen i stället för att lämna
över den ser du i körningen i \cref{ex:events2}, och att hälsningen går
att köra utan fönster ser du i \cref{ex:check-greeting}.  Några
påståenden handlar om vad \mintinline{python}{tkinter} och Tk
\emph{gör} --- att \mintinline{python}{mainloop} kör tills fönstret
stängs, att en bindning anropar sin funktion med en händelse, att en
linje får sin färg med \texttt{fill} --- och för dem är bibliotekens egen
dokumentation den rätta källan; adresserna står i fotnoterna där
påståendena görs, med datum.

Men några påståenden kan ingen körning avgöra:
\begin{enumerate}
  \item att nybörjare förväxlar att koppla en funktion till en händelse
    med att köra den, och är osäkra på vem som anropar den och när
    (\cref{sec:modell})\ifbool{ltnote}{, och att
    \mintinline{python}{mainloop} kan läsas som \enquote{visa fönstret}
    (lärarnoten i \cref{sec:forsta-fonster})}{};
  \item att det är en etablerad princip att hålla gränssnitt och
    programlogik isär, och att det gör logiken lättare att pröva
    (\cref{sec:utan-fonster});
  % The premise is used only in the teacher's notes.
  \ifbool{ltnote}{%
  \item att det som varierar mot en invariant bakgrund är det som kan
    urskiljas, den variationsteoretiska premiss som lärarnoterna i
    marginalen bygger på;
  }{}%
  \item att man bör ärva när den nya klassen är en variant av den gamla
    och föredra komposition när den har en, och att regeln kommer från
    \emph{Design Patterns} (\cref{sec:klassen}).
\end{enumerate}
Hur vet man om sådant stämmer?  Samma fråga möter dig varje gång du läser
en lärobok, en webbsida eller ett svar från en språkmodell --- och varje
gång du själv ska skriva en rapport.  Arbetsgången är densamma i alla
kursens föreläsningar och övningar och beskrivs i föreläsningen
\emph{Algoritmiskt tänkande}, bilaga~A: hur ett påstående görs till en
fråga som kan besvaras, hur källan väljs efter frågan och söks fram
ämnesdrivet och tvåsidigt, med en motsökning efter invändningar, hur
fulltexten läses och spåret skrivs ner, och hur slutsatsen dras med sina
förbehåll.  Påståendet om nybörjare prövas i
\cref{app:grafik-handelser}, principen om gränssnitt och logik i
\cref{app:grafik-separation}, och tumregeln om arv är redan prövad i
föreläsningen \emph{Klasser och objekt}.  \Cref{tab:grafik-pastaenden}
sammanfattar frågorna och svaren.

\begin{table}[htbp]
  \centering
  \caption{Föreläsningens påståenden som frågor, och var svaren finns.
    De två första besvaras i den här föreläsningens bilagor, det tredje i
    en bilaga till en tidigare föreläsning.\ifbool{ltnote}{  Det fjärde
    är en premiss som hela kursen delar.}{}}
  \label{tab:grafik-pastaenden}
  \begin{tabular}{@{}p{0.48\linewidth}p{0.42\linewidth}@{}}
    \toprule
    Fråga & Svar \\
    \midrule
    Har nybörjare svårt med händelsestyrda program och
    \foreignlanguage{english}{callback}-funktioner, och vilka
    missuppfattningar är belagda?
      & Delvis: osäkerheten om vem som anropar en funktion och när, och
        förväxlingen mellan att koppla och att köra den, är belagda (i
        JavaScript); att ämnet i sig är svårt är inte belagt:
        \cref{app:grafik-handelser} \\
    Är det en etablerad princip att hålla gränssnitt och programlogik
    isär, och gör det programmet lättare att pröva?
      & Ja på den första; på den andra med förbehåll, eftersom nyttan
        bara är belagd indirekt: \cref{app:grafik-separation} \\
    Bör man föredra komposition framför arv, och kommer regeln från
    \emph{Design Patterns}?
      & Ja på båda, som en tumregel med förbehåll: belagt i
        \emph{Klasser och objekt}, bilaga~D \\
    % A premise only the teacher's notes use.
    \ifbool{ltnote}{%
    Är det som varierar mot en invariant bakgrund det som kan urskiljas?
      & Premiss från variationsteorin, med källor i lärarnoternas
        källfil; samma i alla kursens föreläsningar \\
    }{}%
    \bottomrule
  \end{tabular}
\end{table}


% The results sections below quote many sources, and every quote carries
% its citation as a margin note; at the notes' normal size the notes of
% the densest pages were pushed onto the next page.  From here on they
% are set one size smaller, as the course book sets them everywhere.
\renewcommand{\foottextfont}{\scriptsize\mpjustification}


% The results sections below quote many sources, and every quote carries
% its citation as a margin note; at the notes' normal size the notes of
% the densest pages were pushed onto the next page.  From here on they
% are set one size smaller, as the course book sets them everywhere.
\renewcommand{\foottextfont}{\scriptsize\mpjustification}

\chapter{Har nybörjare svårt med händelsestyrda program och
  callback-funktioner?}
\label{app:grafik-handelser}
\chapterprecis{Författaren har ännu inte granskat resultaten i den här
  bilagan i sin helhet.}

I \cref{sec:handelsestyrda} jämför materialet ett program som frågar
med \mintinline{python}{input()} och ett fönster där användaren
bestämmer ordningen, i \cref{sec:modell} bygger det en egen
händelseslinga (\foreignlanguage{english}{event loop}) i terminalen och
visar att \mintinline{python}{greet} och \mintinline{python}{greet()}
är två olika saker, och i \cref{sec:forsta-fonster} skriver det ut en
rad efter \mintinline{python}{mainloop()} för att visa när den körs.
Efter körningen i \cref{sec:modell} påstår materialet att förväxlingen
är känd från studier av nybörjare: att en del tror att en funktion körs
redan när den kopplas till en händelse, och att flera är osäkra på vem
som anropar den och när\autocite{Lehtinen2021Struggle,%
Lukkarinen2021Listener}.  Att lämna över en funktion, en
återanropsfunktion (\foreignlanguage{english}{callback}), i stället för
att anropa den är samma skillnad i tkinter som mellan
\mintinline{python}{command=greet} och
\mintinline{python}{command=greet()}.  Frågan som ska besvaras är:
\emph{har nybörjare svårt med händelsestyrda program och
återanropsfunktioner, och vilka missuppfattningar är belagda?}  Frågan
har tre delar: om händelsestyrd programmering
(\foreignlanguage{english}{event-driven programming}) alls är belagd
som svår för nybörjare, om missuppfattningar om hur ett händelsestyrt
program styrs är belagda --- att programmet inte längre går uppifrån och
ned utan väntar på händelser, att dess funktioner anropas av
händelseslingan och att till exempel koden efter
\mintinline{python}{mainloop()} inte körs förrän fönstret stängs ---
och om förväxlingen mellan att lämna över och att anropa en funktion
är det.  Den första delen prövas eftersom påståendet ofta görs i
allmän form, även om materialet självt inte gör det.
\Cref{sec:handelser-metod-grafik} beskriver hur
påståendet prövats, \cref{sec:handelser-resultat-grafik} vad som fanns,
\cref{sec:handelser-diskussion-grafik} vad som begränsar svaret,
\cref{sec:handelser-slutsats-grafik} ger det och
\cref{sec:handelser-fortsatt-grafik} säger vad som återstår att göra;
\cref{tab:sok-grafik-handelser} sist i kapitlet listar de träffar som
rör påståendet.

\section{Metod}
\label{sec:handelser-metod-grafik}

Frågan söktes ämnesdrivet och tvåsidigt den 30 september 2026 med
verktyget \texttt{scholar}, mot de databaser som svarade: OpenAlex, IEEE
Xplore, Web of Science och Scopus.  DBLP svarade hela dagen med en sida
som skulle skilja människor från program i stället för med träffar, och
Semantic Scholar avvisade sin nyckel; båda saknas.  Stödfrågorna
(S1--S10 i \cref{tab:sok-queries-grafik-handelser}) söker svårigheter
och missuppfattningar hos nybörjare i händelsestyrd programmering och
programmering av grafiska gränssnitt, i hur ett händelsestyrt program
styrs, i samtidighet i blockspråk som Scratch och Alice, i funktioner
som värden, återanropsfunktioner och högre ordningens funktioner, i
funktionsanrop och i asynkron programmering.  Motfrågorna (M1--M10 i
\cref{tab:sok-queries-M-grafik-handelser})
använder samma begrepp men söker studier där händelsestyrd
programmering eller funktioner som värden visat sig lätta, naturliga
eller inte svårare, kurser som börjar med händelser, och studier där
svårigheterna förklaras av verktygen eller undervisningens ordning.
Tre motfrågor lägger till fältets egna namn för närliggande begrepp:
\enquote{natural programming} och \enquote{end-user programming} för hur
personer utan programmeringsvana själva uttrycker händelser (M6, M10),
och \enquote{events-first} för kurser som börjar med händelser (M1, M5).

% Author's note (verification and reproduction): see the block above the
% hit-list \input at the end of this chapter.  The query table is
% transcribed from the session's per-provider hit counts (matrix.tsv in
% the research directory, one row per query and provider).
\begin{table}[htbp]
  \centering
  \caption{Stödfrågor och antal träffar per databas, händelsestyrd
    programmering och återanropsfunktioner, 30 september 2026, verktyget
    \texttt{scholar}.  S = stödfråga; motfrågorna står i
    \cref{tab:sok-queries-M-grafik-handelser}.  Databaser: OpenAlex
    (OA; inte \foreignlanguage{english}{open access}), IEEE Xplore
    (IEEE), Web of Science (WoS), Scopus och DBLP.  Högst 40 träffar
    hämtades per databas och fråga, så 40 betyder \enquote{minst 40}.
    Frågan visas i den booleska form som skickades till OpenAlex och IEEE
    Xplore; Web of Science fick samma sträng som
    \texttt{TS=(\dots)} och Scopus som \texttt{TITLE-ABS-KEY(\dots)}.
    DBLP svarade inte (--): tjänsten gav en kontrollsida i stället för
    träffar hela dagen.  Semantic Scholar avvisade sin nyckel och saknas
    helt.  Frågorna står i den ordning de ställdes.}
  \label{tab:sok-queries-grafik-handelser}
  \scriptsize
  \begin{tabular}{@{}lp{0.46\linewidth}rrrrr@{}}
    \toprule
    & Nyckelord & OA & IEEE & WoS & Scopus & DBLP \\
    \midrule
    S1 & (\enquote{event-driven} \textsc{or} \enquote{event driven}
       \textsc{or} \enquote{event-based} \textsc{or} \enquote{graphical
       user interface} \textsc{or} GUI) \textsc{and} (novice \textsc{or}
       novices \textsc{or} beginners \textsc{or} CS1 \textsc{or}
       \enquote{introductory programming}) \textsc{and} (difficulties
       \textsc{or} difficult \textsc{or} misconceptions \textsc{or}
       \enquote{mental models})
       & 40 & 39 & 40 & 40 & -- \\
    S2 & (\enquote{event-driven} \textsc{or} \enquote{event driven}
       \textsc{or} \enquote{event loop} \textsc{or} \enquote{event
       handler} \textsc{or} \enquote{event handlers} \textsc{or}
       \enquote{event handling} \textsc{or} callback \textsc{or}
       callbacks) \textsc{and} (\enquote{control flow} \textsc{or}
       \enquote{flow of control} \textsc{or} \enquote{sequential
       execution} \textsc{or} \enquote{inversion of control})
       \textsc{and} (students \textsc{or} novices \textsc{or} beginners
       \textsc{or} learners)
       & 40 & 4 & 4 & 6 & -- \\
    S3 & (callback \textsc{or} callbacks \textsc{or} \enquote{higher-order
       function} \textsc{or} \enquote{higher-order functions} \textsc{or}
       \enquote{first-class functions} \textsc{or} \enquote{functions as
       values} \textsc{or} \enquote{function reference} \textsc{or}
       lambda) \textsc{and} (novice \textsc{or} novices \textsc{or}
       students \textsc{or} beginners) \textsc{and} (misconception
       \textsc{or} misconceptions \textsc{or} difficulties \textsc{or}
       errors \textsc{or} mistakes)
       & 40 & 22 & 40 & 40 & -- \\
    S4 & (Scratch \textsc{or} Alice \textsc{or} \enquote{block-based}
       \textsc{or} \enquote{App Inventor}) \textsc{and}
       (\enquote{event-driven} \textsc{or} events \textsc{or}
       \enquote{event handling} \textsc{or} concurrency \textsc{or}
       parallelism \textsc{or} synchronization) \textsc{and}
       (misconception \textsc{or} misconceptions \textsc{or} difficulties
       \textsc{or} \enquote{bad habits} \textsc{or} errors)
       & 39 & 40 & 40 & 39 & -- \\
    S5 & (\enquote{function call} \textsc{or} \enquote{function calls}
       \textsc{or} \enquote{function calling} \textsc{or}
       \enquote{procedure call}) \textsc{and} (novice \textsc{or} novices
       \textsc{or} beginners \textsc{or} CS1) \textsc{and} (misconception
       \textsc{or} misconceptions \textsc{or} \enquote{mental model}
       \textsc{or} \enquote{mental models})
       & 40 & 2 & 0 & 1 & -- \\
    S6 & (\enquote{event-driven programming} \textsc{or}
       \enquote{event-driven programs} \textsc{or} \enquote{GUI
       programming} \textsc{or} \enquote{event handling} \textsc{or}
       \enquote{event handlers}) \textsc{and} (teaching \textsc{or}
       learning \textsc{or} students \textsc{or} novices \textsc{or} CS1
       \textsc{or} \enquote{introductory programming})
       & 40 & 40 & 40 & 40 & -- \\
    S7 & (\enquote{event-driven} \textsc{or} events \textsc{or}
       concurrency \textsc{or} concurrent \textsc{or} parallelism
       \textsc{or} callbacks) \textsc{and} (novices \textsc{or} novice
       \textsc{or} \enquote{high school} \textsc{or} children \textsc{or}
       CS1) \textsc{and} (misconceptions \textsc{or} \enquote{mental
       models} \textsc{or} \enquote{mental model} \textsc{or}
       \enquote{notional machine})
       & 40 & 27 & 40 & 40 & -- \\
    S8 & (\enquote{higher-order functions} \textsc{or}
       \enquote{higher-order function} \textsc{or} \enquote{first-class
       functions} \textsc{or} \enquote{functions as values} \textsc{or}
       \enquote{functions as arguments} \textsc{or} callbacks \textsc{or}
       \enquote{lambda expressions}) \textsc{and} (learning \textsc{or}
       teaching \textsc{or} students \textsc{or} novices) \textsc{and}
       (understanding \textsc{or} difficulties \textsc{or} difficulty
       \textsc{or} comprehension)
       & 40 & 31 & 40 & 40 & -- \\
    S9 & (\enquote{missing parentheses} \textsc{or} \enquote{without
       parentheses} \textsc{or} \enquote{function object} \textsc{or}
       \enquote{function objects} \textsc{or} \enquote{function
       reference} \textsc{or} \enquote{method reference} \textsc{or}
       \enquote{passing a function} \textsc{or} \enquote{passing
       functions}) \textsc{and} (novice \textsc{or} novices \textsc{or}
       students \textsc{or} beginners \textsc{or} CS1) \textsc{and}
       (errors \textsc{or} mistakes \textsc{or} misconception \textsc{or}
       misconceptions \textsc{or} difficulties)
       & 40 & 16 & 1 & 6 & -- \\
    S10 & (asynchronous \textsc{or} async \textsc{or} callbacks
       \textsc{or} callback \textsc{or} promises \textsc{or}
       \enquote{event loop}) \textsc{and} (JavaScript \textsc{or} Python
       \textsc{or} programming) \textsc{and} (novices \textsc{or} novice
       \textsc{or} students \textsc{or} learners) \textsc{and}
       (misconceptions \textsc{or} difficulties \textsc{or}
       \enquote{mental model} \textsc{or} \enquote{mental models})
       & 40 & 40 & 40 & 39 & -- \\
    \bottomrule
  \end{tabular}
\end{table}

\begin{table}[htbp]
  \centering
  \caption{Motfrågor och antal träffar per databas, händelsestyrd
    programmering och återanropsfunktioner, 30 september 2026, verktyget
    \texttt{scholar}.  M = motfråga; förkortningar, tak och
    frågeformer som i \cref{tab:sok-queries-grafik-handelser}.  DBLP
    svarade inte (--) och Semantic Scholar saknas.}
  \label{tab:sok-queries-M-grafik-handelser}
  \scriptsize
  \begin{tabular}{@{}lp{0.46\linewidth}rrrrr@{}}
    \toprule
    & Nyckelord & OA & IEEE & WoS & Scopus & DBLP \\
    \midrule
    M1 &(\enquote{event-driven} \textsc{or} \enquote{event driven}
       \textsc{or} \enquote{event-based} \textsc{or} \enquote{event loop}
       \textsc{or} \enquote{event handler} \textsc{or} \enquote{event
       handling} \textsc{or} \enquote{graphical user interface}
       \textsc{or} GUI \textsc{or} \enquote{events-first} \textsc{or}
       \enquote{events first}) \textsc{and} (novice \textsc{or} novices
       \textsc{or} beginners \textsc{or} CS1 \textsc{or}
       \enquote{introductory programming}) \textsc{and} (easier
       \textsc{or} easy \textsc{or} natural \textsc{or} intuitive
       \textsc{or} \enquote{no difference} \textsc{or} \enquote{no
       significant} \textsc{or} motivating \textsc{or} engaging)
       & 40 & 40 & 40 & 40 & -- \\
    M2 & (callback \textsc{or} callbacks \textsc{or} \enquote{higher-order
       function} \textsc{or} \enquote{higher-order functions} \textsc{or}
       \enquote{first-class functions} \textsc{or} \enquote{functions as
       values} \textsc{or} \enquote{function reference} \textsc{or}
       lambda \textsc{or} \enquote{function call} \textsc{or}
       \enquote{function calls}) \textsc{and} (novice \textsc{or} novices
       \textsc{or} students \textsc{or} beginners) \textsc{and} (easier
       \textsc{or} easy \textsc{or} natural \textsc{or} intuitive
       \textsc{or} \enquote{no difference} \textsc{or} \enquote{no
       significant} \textsc{or} \enquote{not difficult})
       & 40 & 40 & 40 & 40 & -- \\
    M3 & (Scratch \textsc{or} Alice \textsc{or} \enquote{block-based}
       \textsc{or} \enquote{App Inventor}) \textsc{and}
       (\enquote{event-driven} \textsc{or} events \textsc{or}
       \enquote{event handling} \textsc{or} concurrency \textsc{or}
       parallelism) \textsc{and} (novice \textsc{or} novices \textsc{or}
       children \textsc{or} students) \textsc{and} (easier \textsc{or}
       easy \textsc{or} natural \textsc{or} intuitive \textsc{or}
       successful \textsc{or} \enquote{no difference})
       & 40 & 21 & 19 & 32 & -- \\
    M4 & (\enquote{event-driven} \textsc{or} \enquote{event driven}
       \textsc{or} \enquote{event handling} \textsc{or} GUI \textsc{or}
       \enquote{graphical user interface} \textsc{or} callbacks
       \textsc{or} \enquote{higher-order functions}) \textsc{and} (novice
       \textsc{or} novices \textsc{or} students \textsc{or} CS1)
       \textsc{and} (misconceptions \textsc{or} difficulties) \textsc{and}
       (replication \textsc{or} contradictory \textsc{or} \enquote{not
       supported} \textsc{or} limitations \textsc{or} \enquote{teaching
       order} \textsc{or} \enquote{order of topics} \textsc{or}
       \enquote{events-first} \textsc{or} \enquote{objects-first})
       & 40 & 2 & 3 & 5 & -- \\
    M5 & (\enquote{event-driven programming} \textsc{or}
       \enquote{event-driven programs} \textsc{or} \enquote{GUI
       programming} \textsc{or} \enquote{event handling} \textsc{or}
       \enquote{event handlers} \textsc{or} \enquote{events-first})
       \textsc{and} (teaching \textsc{or} learning \textsc{or} students
       \textsc{or} novices \textsc{or} CS1 \textsc{or}
       \enquote{introductory programming}) \textsc{and} (easier
       \textsc{or} easy \textsc{or} natural \textsc{or} \enquote{simple
       enough} \textsc{or} successful \textsc{or} \enquote{no difference}
       \textsc{or} \enquote{no significant} \textsc{or} motivating)
       & 40 & 23 & 24 & 40 & -- \\
    M6 & (\enquote{non-programmers} \textsc{or} \enquote{natural
       programming} \textsc{or} children \textsc{or} novices \textsc{or}
       \enquote{high school} \textsc{or} CS1) \textsc{and}
       (\enquote{event-based} \textsc{or} \enquote{event-driven}
       \textsc{or} \enquote{rule-based} \textsc{or} events \textsc{or}
       concurrency \textsc{or} concurrent \textsc{or} callbacks)
       \textsc{and} (natural \textsc{or} intuitive \textsc{or}
       commonsense \textsc{or} \enquote{prior knowledge} \textsc{or}
       spontaneous \textsc{or} \enquote{no misconceptions})
       & 40 & 40 & 40 & 40 & -- \\
    M7 & (\enquote{higher-order functions} \textsc{or}
       \enquote{higher-order function} \textsc{or} \enquote{first-class
       functions} \textsc{or} \enquote{functions as values} \textsc{or}
       \enquote{functions as arguments} \textsc{or} callbacks \textsc{or}
       \enquote{lambda expressions}) \textsc{and} (learning \textsc{or}
       teaching \textsc{or} students \textsc{or} novices) \textsc{and}
       (easier \textsc{or} easy \textsc{or} natural \textsc{or} intuitive
       \textsc{or} \enquote{no difference} \textsc{or} \enquote{no
       significant} \textsc{or} successful)
       & 40 & 26 & 33 & 40 & -- \\
    M8 & (\enquote{missing parentheses} \textsc{or} \enquote{without
       parentheses} \textsc{or} \enquote{function object} \textsc{or}
       \enquote{function objects} \textsc{or} \enquote{function
       reference} \textsc{or} \enquote{method reference} \textsc{or}
       \enquote{passing a function} \textsc{or} \enquote{passing
       functions}) \textsc{and} (novice \textsc{or} novices \textsc{or}
       students \textsc{or} beginners \textsc{or} CS1) \textsc{and}
       (easier \textsc{or} easy \textsc{or} natural \textsc{or} intuitive
       \textsc{or} \enquote{no difference} \textsc{or} \enquote{no
       significant} \textsc{or} successful)
       & 40 & 13 & 2 & 8 & -- \\
    M9 & (asynchronous \textsc{or} async \textsc{or} callbacks \textsc{or}
       callback \textsc{or} promises \textsc{or} \enquote{event loop})
       \textsc{and} (JavaScript \textsc{or} Python \textsc{or}
       programming) \textsc{and} (novices \textsc{or} novice \textsc{or}
       students \textsc{or} learners) \textsc{and} (easier \textsc{or}
       easy \textsc{or} natural \textsc{or} intuitive \textsc{or}
       \enquote{no difference} \textsc{or} \enquote{no significant}
       \textsc{or} successful)
       & 40 & 40 & 40 & 40 & -- \\
    M10 & (\enquote{natural programming} \textsc{or}
       \enquote{non-programmers} \textsc{or} nonprogrammers \textsc{or}
       \enquote{end-user programmers} \textsc{or} \enquote{end-user
       programming}) \textsc{and} (\enquote{event-based} \textsc{or}
       \enquote{event-driven} \textsc{or} \enquote{rule-based} \textsc{or}
       events \textsc{or} \enquote{event handling} \textsc{or} callbacks)
       & 40 & 34 & 40 & 40 & -- \\
    \bottomrule
  \end{tabular}
\end{table}

Sammanlagt lämnade databaserna 2\,521 poster.  Flera databaser lämnar
samma arbete, och två poster räknas som ett arbete när titlarna är
desamma bortsett från versaler och skiljetecken; efter den
sammanslagningen återstår 1\,812 arbeten, och det är arbeten som räknas
här och i \cref{tab:sok-grafik-handelser}.  Alla klassades: först med
en språkmodell mot en beskrivning av frågan och fyra kategorier ---
stöder påståendet, kvalificerar eller motsäger det, angränsande
delämne, annat ämne --- och sedan för hand: varje post som modellen
satt i någon av de två första kategorierna, och varje post med
konfidens under 0,7, lästes i sammanfattning.  Av de 80 poster som
lästes så behöll 24 modellens beslut; 54 flyttades från stödjande
eller kvalificerande till angränsande eller annat ämne, oftast
beskrivningar av ett verktyg eller en kurs utan något resultat om
hur man lär sig, och två flyttades åt andra hållet.  Därefter återstod
1\,218 felträffar och 546 angränsande arbeten.  De flesta felträffarna
kom från de breda frågorna S1 och M1, där \enquote{GUI} och
\enquote{graphical user interface} ger all vetenskaplig programvara
med ett grafiskt gränssnitt; det är skälet till att frågorna S6 och M5
ställdes kring själva termen \enquote{event-driven programming}.

\Cref{tab:sok-grafik-handelser} rättades för hand vid genomläsningen,
och rättelserna redovisas därför att en rättelse man inte kan
kontrollera inte är någon rättelse.  Två titlar hade citattecken som
kom ut som obearbetade styrsekvenser, och en titel hade \enquote{C++}
utskrivet som \enquote{C plus plus}.  Två arbeten stod på två rader
var, eftersom databaserna gav dem något olika titlar, och står nu på
var sin rad: Bruce, Danyluk och Murtaghs artikel, och Christensen och
Caspersens, vars undertitel bara en av databaserna hade med.

Ett arbete hämtades utanför sökningen.  Lukkarinens
doktorsavhandling \parencite{Lukkarinen2022Thesis} är ingen
tidskriftsartikel och indexeras inte av databaserna; den hittades när
en av studierna i \cref{sec:handelser-resultat-grafik} söktes upp i
Aalto-universitetets arkiv.  Fulltexterna till Lukkarinen
med fleras två artiklar \parencite{Lukkarinen2021Mapping,
Lukkarinen2021Listener} och till Lu och Krishnamurthis artikel
\parencite{LuKrishnamurthi2024} hämtades från författarnas egna kopior,
eftersom förlagets bibliotek avvisar automatisk hämtning.  Sju av de
21 källor som kapitlet citerar lästes i fulltext; de övriga 14 lästes
bara i sammanfattning, och \cref{sec:handelser-diskussion-grafik}
återkommer till det.  Hur varje källa hittades och verifierades är
antecknat på det sätt som bilagan om arbetsgången beskriver.

\section{Resultat}
\label{sec:handelser-resultat-grafik}

\paragraph{Är händelsestyrd programmering svår för nybörjare?}  Den
enda systematiska genomgången av forskningen om att lära ut och lära
sig händelsestyrd programmering, \textcite{Lukkarinen2021Mapping}, gick
igenom 858 träffar och tog med 105 arbeten fram till hösten 2018.
Den konstaterar att undervisningen har
\textcquote{Lukkarinen2021Mapping}{revealed many challenges related to
teaching and learning edp}, och att flera författare beskriver
\textcquote{Lukkarinen2021Mapping}{the concept as a whole as
challenging or hard to understand}.  En av de äldre studier som
genomgången återger ligger nära det materialet tar upp om
\mintinline{python}{mainloop()}: lärare som undervisat i Pascal hade
svårt att hitta \textcquote{Lukkarinen2021Mapping}{the main program and
the place where the program would finally end, as well as the control
flows of programs developed using Delphi}.  Två studier som sökningen
gav bekräftar bilden med egna data.  \Textcite{Shackleton1998} följde
nybörjare som lärde sig Visual Basic och fann att
\textcquote[193]{Shackleton1998}{many have great difficulty in the
early stages understanding the concepts of event-driven systems}; de
klarade att rita upp fönstret men hade svårt att föreställa sig hur
programmet sedan körs.  \Textcite{WiedenbeckEngebretson2004} lät lärare
tänka högt medan de läste ett händelsestyrt program och fann svårigheter
\textcquote{WiedenbeckEngebretson2004}{comprehending the event-driven
application, given the distributed nature of the code}.  Att
svårigheten tas för given syns också i verktygen: \textcite{Mackin2013}
föreslår ett grafikbibliotek utan händelser, med motiveringen att
\textcquote{Mackin2013}{the idea of threads and event handling can be
difficult to grasp}, men mäter inte den svårigheten.

\paragraph{Missuppfattningar om hur programmet styrs.}  Fram till 2018
fanns ingen studie av sådana missuppfattningar:
\textcquote{Lukkarinen2021Mapping}{We did not find any work related to
students' misconceptions of edp concepts, neither any systematic
qualitative studies of their conceptions}.  Den första som finns är
\textcite{Lukkarinen2021Listener}, som lät vuxna studenter i en
nätkurs i webbprogrammering med JavaScript --- de flesta nybörjare ---
förklara begreppen och en kort kodsnutt med en knapp.  Studenterna
beskrev lyssnaren som något som aktivt väntar på händelser och anropar
hanteraren, och några gav själva registreringen egenskaper den inte har:
\textcquote[8]{Lukkarinen2021Listener}{Two of them argued that calling
addEventListener() waits for or detects clicks. [\dots] A fourth one
more specifically described blocking behaviour in claiming that the
function continues until the button is clicked}.  Författarnas slutsats
är att studenterna var osäkra på
\textcquote[9]{Lukkarinen2021Listener}{exactly who is calling the
listeners, when, how the program execution proceeds after the code of a
listener has been executed}.  Det är samma fråga som
\mintinline{python}{mainloop()} väcker, fast omvänd: studenterna trodde
att ett anrop som inte väntar väntade, medan den som skriver kod efter
\mintinline{python}{mainloop()} tror att ett anrop som väntar inte gör
det.  Lukkarinens avhandling, som sammanfattar gruppens studier, lägger
till att en knapp för en del studenter själv utförde det som händer när
man trycker på den,
\textcquote[45]{Lukkarinen2022Thesis}{instead of just triggering the
execution of some program code}, och konstaterar att den vanligaste
katalogen över missuppfattningar i nybörjarprogrammering inte tar upp
händelsestyrd programmering alls
\parencite[19]{Lukkarinen2022Thesis}.  I blockspråk, där flera skript
körs samtidigt som svar på händelser, fann
\textcite{MeerbaumSalant2013} att högstadieelever lärde sig viktiga
begrepp men hade problem med bland annat samtidighet
(\foreignlanguage{english}{concurrency}), och \textcite{Kolikant2001}
fann att gymnasieelever före en kurs i samtidighet bland annat
\textcquote{Kolikant2001}{attribute parallelism where it does not
exist}.  Den särskilda missuppfattning som materialet beskriver ---
att \mintinline{python}{mainloop()} bara visar fönstret och att koden
efter den körs direkt --- fanns inte beskriven i någon studie.

\paragraph{Att lämna över en funktion eller att anropa den.}
\Textcite{Lehtinen2021Struggle} ställde frågor om studenternas egen
kod i samma nätkurs som ovan.  Sju svar (21~procent) på frågorna om
händelsehanteringen tydde på förvirring om
\textcquote{Lehtinen2021Struggle}{what it means to attach (to add, to
register) an event handler, when they are actually executed, and who
actually calls them}; ett svar förklarade att registreringen
\enquote{executes the playNote event handler}, vilket författarna läser
som \textcquote{Lehtinen2021Struggle}{the incorrect idea of
addEventListener() method calling the given event handler}.  Det är
begreppsligt samma förväxling som mellan
\mintinline{python}{command=f} och \mintinline{python}{command=f()}:
att koppla en funktion till en händelse tas för att köra den.
Förväxlingen finns också utanför händelser.  \Textcite{LuKrishnamurthi2024}
kallar den missuppfattningen att
\textcquote[106:17]{LuKrishnamurthi2024}{Functions are not considered
first-class values. They can't be bound to other variables, passed as
arguments, or referred to by data structures}, och elva procent av
svaren på ett program som lämnar över en funktion som argument visade
den --- hos studenter i årskurs tre och fyra i en valbar kurs om
programspråk, alltså inte nybörjare.  Hos nybörjare i Python finns
den spegelvända förväxlingen, ett funktionsnamn utan anropsparenteser:
\textcite[69]{Bosse2021} katalogiserar \enquote{Missing parentheses in
\mbox{\enquote{input}}} bland misstagen hos 166 studenter i en första
programmeringskurs.  I ett kontrollerat försök skrev studenter
långsammare korrekta program med anonyma funktioner i C\texttt{++} än
utan \parencite{Uesbeck2016}.  Ingen studie av just
\mintinline{python}{command=f()} i Python eller tkinter kom fram.

\paragraph{Motsökningen.}  Motsökningen gav två slags resultat, och
båda kvalificerar påståendet.  Det första är kurser där händelsestyrd
programmering från början har gått bra.  \Textcite{BruceDanyluk2001}
byggde en introduktionskurs i Java kring händelser från första veckan
och skriver att deras erfarenhet \textcquote{BruceDanyluk2001}{runs
counter to the prevailing sense that these techniques would add
complexity to the content of CS 1}; samma grupp menar senare att
händelserna till och med gör kursens övriga innehåll lättare
\parencite{BruceDanyluk2004}.  \Textcite{ChristensenCaspersen2002}
introducerade omvänd styrning (\foreignlanguage{english}{inversion of
control}) i en första kurs genom ett litet ramverk, så att studenterna
inte överväldigades av detaljerna i Javas stora
gränssnittsbibliotek.  \Textcite{Werner2012} fann att många av 231
spel som högstadieelever byggt i Alice använde händelsehanterare och
samtidighet på ett fungerande sätt, och \textcite{Wang2016} att elva
barn i årskurs ett och två fick en första förståelse av
händelsehantering med programmeringsblock att ta i.  Från andra
hållet byggde \textcite{Pane2002}, utifrån studier av hur personer utan
programmeringsvana beskriver lösningar, ett händelsebaserat programspråk
för barn med operationer som \textcquote{Pane2002}{match the way
non-programmers express problem solutions}.  I det enda
kontrollerade försöket som kom fram förstod nybörjare ett litet
händelsestyrt program ungefär lika bra som ett likvärdigt
objektorienterat, och frågorna om programflödet besvarades väl i båda
\parencite{KhazaeiJackson2002}.  Det andra resultatet står i
genomgången själv: \textcquote{Lukkarinen2021Mapping}{despite a number
of personal reflections on the challenge of learning edp, there was no
consensus that learning the topic per se is difficult. More often the
challenges were related to the actual tools available.}  För funktioner
som värden gäller detsamma: \textcite{KrishnamurthiFisler2021} fann att
studenter först hade flera svårigheter med högre ordningens funktioner
men senare klarade dem \textcquote{KrishnamurthiFisler2021}{quite
well}, och \textcite{Rivera2022} att studenter tidigt i sin utbildning
kände igen enskilda sådana funktioner och kunde sätta ihop dem till
fungerande program.  Ingen studie som hävdar att nybörjare saknar
svårigheter med händelsestyrda program kom fram.

\section{Diskussion och begränsningar}
\label{sec:handelser-diskussion-grafik}

Fyra saker begränsar svaret.  Den första är att ingen källa gäller
Python eller tkinter: de direkta beläggen för missuppfattningar kommer
från JavaScript \parencite{Lukkarinen2021Listener,Lehtinen2021Struggle},
de äldre från Visual Basic och Delphi, blockspråken från Scratch och
Alice, och förväxlingen mellan funktion och anrop från ett språk inbyggt
i Racket \parencite{LuKrishnamurthi2024}.  Begreppen --- vem som anropar
hanteraren, vilket anrop som väntar, skillnaden mellan en funktion och
dess anrop --- är desamma i Python, så överföringen är rimlig, men den är
inte prövad.  Den andra är populationerna: de två JavaScript-studierna
gäller vuxna i en öppen nätkurs, med 24 och 26 analyserade svar, och
författarna varnar själva för att frekvenserna
\textcquote[9]{Lukkarinen2021Listener}{should be considered describing
this data set only}; siffran elva procent gäller studenter i årskurs tre
och fyra, inte nybörjare.  Hur vanliga missuppfattningarna är hos våra
studenter säger alltså ingen källa.  Den tredje är verifieringen: sju
källor lästes i fulltext, men de 14 övriga, bland dem hela motsidan
utom genomgången, kunde bara läsas i sammanfattning, eftersom ACM:s och
IEEE:s bibliotek avvisar automatisk hämtning och källorna inte fanns i
författarens bibliotek på läsplattan; det står i källfilen.  Shackletons
studie är dessutom en masteruppsats och Mackins en konferenstext utan
mätning.  Den fjärde är täckningen: DBLP och Semantic
Scholar saknas, träffantalen i
\cref{tab:sok-queries-grafik-handelser,tab:sok-queries-M-grafik-handelser}
är tak på 40, OpenAlex relevansordning begravde de få relevanta
arbetena bland felträffar, och klassningen är modellstödd: de 13
stödjande och 15 kvalificerande posterna i
\cref{tab:sok-grafik-handelser} som inte citeras är bara lästa i
sammanfattning eller, där en sådan saknades, bedömda på titeln, så
antalen är övre gränser.

\section{Slutsats}
\label{sec:handelser-slutsats-grafik}

Delvis.  Att nybörjare missförstår hur ett händelsestyrt program styrs
är belagt: de är osäkra på vem som anropar en händelsehanterare, när det
sker och vilket anrop som väntar
\parencite{Lukkarinen2021Listener,Lukkarinen2022Thesis}, och en del tar
det att koppla en hanterare till en händelse för att köra den
\parencite{Lehtinen2021Struggle} --- samma förväxling som mellan
\mintinline{python}{command=f} och \mintinline{python}{command=f()}.
Att funktioner inte behandlas som värden är belagt även hos erfarna
studenter \parencite{LuKrishnamurthi2024}.  Men påståendet att
händelsestyrd programmering \emph{är svår} för nybörjare är starkare än
underlaget: den enda systematiska genomgången fann ingen enighet om det
och lade det mesta av svårigheten på verktygen
\parencite{Lukkarinen2021Mapping}, och kurser som börjat med händelser
rapporterar att det gick bra
\parencite{BruceDanyluk2001,ChristensenCaspersen2002}.  De två
särskilda formerna i materialet --- \mintinline{python}{mainloop()} som
\enquote{visa fönstret} och \mintinline{python}{command=f()} --- är inte
studerade, utan är tillämpningar på tkinter av de belagda
missförstånden.  Påståendet i \cref{sec:modell} --- att en del tror att
en funktion körs redan när den kopplas till en händelse, och att flera
är osäkra på vem som anropar den och när --- håller därför i den
styrka det har, med förbehållet att studierna gäller JavaScript och
vuxna nybörjare i en nätkurs och inte säger hur vanliga missförstånden
är hos våra studenter.  Att området i sig är svårt kan materialet
däremot inte säga, och med stöd av motsidan kan det lägga till att
händelser går att lära ut från början om verktyget är enkelt.

\section{Fortsatt arbete}
\label{sec:handelser-fortsatt-grafik}

Sökningen lämnade tre saker ogjorda.  Två databaser saknas: DBLP
svarade inte på någon fråga och Semantic Scholar avvisade sin nyckel,
så samma 20 frågor bör ställas där när tjänsterna svarar; DBLP indexerar
just de konferenser om programmeringsundervisning där de flesta
arbetena ovan finns.  Hela motsidan är vidare läst bara i
sammanfattning.  Tyngst väger Bruce, Danyluk och Murtaghs rapport
\autocite{BruceDanyluk2001} och Khazaei och Jacksons försök
\autocite{KhazaeiJackson2002}: en läsning i fulltext skulle visa hur
stort försöket var, vilket språk det gällde och om kursen med händelser
från början mätte något alls.  Detsamma gäller de två studierna av
högre ordningens funktioner \autocite{KrishnamurthiFisler2021,
Rivera2022}, vars språk och population sammanfattningarna inte anger.
Den tredje luckan är de studier som genomgången själv refererar men
som sökningen inte gav, främst Halland och Malans studie av lärare som
gick från Pascal till Delphi; den citeras här i andra hand.

För ett säkrare svar saknas tre slags studier.  Ingen av källorna gäller
Python eller tkinter, och ingen mäter hur vanliga missförstånden är
hos studenter i en första programmeringskurs.  Ett diagnostiskt prov i
kursens egen population --- vad händer med raden efter
\mintinline{python}{mainloop()}, vad skickas med
\mintinline{python}{command=f()}, vem anropar funktionen --- skulle ge
den frekvensen.  Ingen studie jämför heller hur man bäst undanröjer
missförstånden; ett försök som jämför materialets upplägg, att först
bygga en egen händelseslinga i terminalen, med att gå direkt till
tkinter skulle visa om modellen bär.  Faller det ut till modellens
fördel bör den stå kvar som materialets första steg; faller det ut åt
andra hållet kan materialet gå direkt till fönstren och i stället
lägga tiden på förväxlingen mellan \mintinline{python}{f} och
\mintinline{python}{f()}, som är den mest direkt belagda.

\section{Träffar som rör påståendet}
\label{sec:handelser-traffar-grafik}

\Cref{tab:sok-grafik-handelser} listar de poster som rör påståendet
--- citerade, stödjande och kvalificerande --- med skälet till varje
beslut; de angränsande och de felträffar som uteslöts anges som antal i
tabelltexten, och antalet träffar per fråga står i
\cref{tab:sok-queries-grafik-handelser,tab:sok-queries-M-grafik-handelser}.
% Author's note (verification and reproduction):
%   scholar session : grafik-handelser
%   exports         : litteratursokning/grafik-handelser.csv,
%                     litteratursokning/grafik-handelser-full.tex
%                     (hand-corrected as the method section says; the
%                     generated original is grafik-handelser-full.generated.tex
%                     in the research directory)
%   searched        : 2026-09-30; providers openalex ieee wos scopus
%                     (dblp: every request answered with a bot-check
%                     page, "non-JSON response (HTTP 200)", also on the
%                     uni-trier mirror, retried at the end -- absent;
%                     s2: key rejected HTTP 403 -- absent); limit 40 per
%                     provider and query; wos as TS=(...), scopus as
%                     TITLE-ABS-KEY(...); exact strings in the session
%                     and in tab:sok-queries-grafik-handelser
%   classification  : scholar llm context + scholar llm classify
%                     --no-examples --no-enrich -n 150 (13 batches),
%                     model openai-codex/gpt-5.6-sol; human review of
%                     every supports/qualifies row and every row with
%                     confidence < 0.7 (80 rows: 24 confirmed, 54 kept ->
%                     discarded, 2 discarded -> kept) via scholar sessions
%                     decide (decide.sh, export.sh in the research dir)
%   theme tags      : edp-difficulty, control-flow, attach-vs-call, hof,
%                     events-first
%   hit-list table  : scholar sessions export grafik-handelser -f table
%                     --bearing-only --lang sv
%                     --label sok-grafik-handelser
%                     --track "händelsestyrda program och
%                     återanropsfunktioner"
%                     --theme edp-difficulty="händelsestyrd programmering
%                     som svårighet" --theme control-flow="hur programmet
%                     styrs" --theme attach-vs-call="lämna över eller
%                     anropa en funktion" --theme hof="funktioner som
%                     värden" --theme events-first="händelser från början"
%                     -o litteratursokning/grafik-handelser-full
%   provenance      : grafik-handelser.bib, keys Lukkarinen2021Mapping
%                     Lukkarinen2021Listener Lehtinen2021Struggle
%                     Lukkarinen2022Thesis Shackleton1998
%                     WiedenbeckEngebretson2004 Mackin2013
%                     MeerbaumSalant2013 Kolikant2001 BruceDanyluk2001
%                     BruceDanyluk2004 KhazaeiJackson2002 Werner2012
%                     ChristensenCaspersen2002 Wang2016 Pane2002
%                     LuKrishnamurthi2024 KrishnamurthiFisler2021
%                     Rivera2022 Uesbeck2016 Bosse2021
%   known items     : Lukkarinen2022Thesis (not indexed; aaltodoc)
%   full text       : Lukkarinen2021Mapping, Lukkarinen2021Listener
%                     (aaltodoc manuscripts), Lehtinen2021Struggle and
%                     Bosse2021 (arXiv), Lukkarinen2022Thesis (aaltodoc),
%                     Shackleton1998 (VU repository), LuKrishnamurthi2024
%                     (author copy, cs.brown.edu)
%   abstract-level  : the other 14 -- ACM and IEEE refuse automated
%                     download, not on the user's reMarkable; for the user
% --- scholar LaTeX fragment --------------------------------------------
% Generated by: scholar sessions export -f table
% Session: grafik-handelser
% \input this file from the body of a document whose
% preamble loads:
%   \usepackage{longtable}
%   \usepackage{booktabs}
%   \usepackage{array}
% ----------------------------------------------------------------------

\begingroup\footnotesize
\begin{longtable}{@{}%
  >{\raggedright\arraybackslash}p{0.44\textwidth}c%
  >{\raggedright\arraybackslash}p{0.13\textwidth}%
  >{\raggedright\arraybackslash}p{0.26\textwidth}@{}}
\caption{Träffar som rör påståendet, händelsestyrda program och återanropsfunktioner (20 citerade, 13 som stöder påståendet och 15 som kvalificerar eller motsäger det, av 1812 unika träffar; de 546 angränsande och 1218 felträffarna listas inte). Skälet till varje inkludering står i sista kolumnen; för maskinklassade rader anges modellens konfidens. Databaser: OpenAlex (OA; inte \foreignlanguage{english}{open access}), IEEE Xplore (IEEE), Web of Science (WoS), Scopus.}\label{tab:sok-grafik-handelser}\\
\toprule
Titel & År & Leverantör & Skäl \\
\midrule
\endfirsthead
\multicolumn{4}{@{}l}{\emph{\tablename~\thetable{} (forts.)}}\\
\toprule
Titel & År & Leverantör & Skäl \\
\midrule
\endhead
\midrule \multicolumn{4}{r@{}}{\emph{forts.\ på nästa sida}}\\
\endfoot
\bottomrule
\endlastfoot
A tangible embedded programming system to convey event-handling concept & 2016 & Scopus & \emph{citerad}: händelser från början \\
An Empirical Study on the Impact of C\texttt{++} Lambdas and Programmer Experience & 2016 & WoS & \emph{citerad}: lämna över eller anropa en funktion \\
An Event Listener or an Event Handler? Students Explain Event-drivenness in JavaScript & 2021 & WoS & \emph{citerad}: hur programmet styrs \\
Catalogs of C and Python Antipatterns by CS1 Students & 2021 & OA & \emph{citerad}: lämna över eller anropa en funktion \\
Children Learning Computer Science Concepts via Alice Game-Programming & 2011 & WoS, Scopus & \emph{citerad}: händelser från början \\
Comprehension Strategies of End-User Programmers in an Event-Driven Application & 2004 & IEEE & \emph{citerad}: händelsestyrd programmering som svårighet \\
Developing Behavioral Concepts of Higher-Order Functions & 2021 & Scopus & \emph{citerad}: funktioner som värden \\
Event driven languages for novice programmers & 1998 & OA & \emph{citerad}: händelsestyrd programmering som svårighet \\
Event-driven programming facilitates learning standard programming concepts & 2004 & Scopus & \emph{citerad}: händelser från början \\
Event-driven Programming in Programming Education: A Mapping Review & 2021 & WoS & \emph{citerad}: händelsestyrd programmering som svårighet \\
Event-driven programming is simple enough for CS1 & 2001 & OA & \emph{citerad}: händelser från början \\
Frameworks in CS1: A different way of introducing event-driven programming & 2002 & OA, Scopus & \emph{citerad}: händelser från början \\
Gardeners and Cinema Tickets: High School Students' Preconceptions of Concurrency & 2001 & OA & \emph{citerad}: hur programmet styrs \\
Identifying and Correcting Programming Language Behavior Misconceptions & 2024 & Scopus & \emph{citerad}: lämna över eller anropa en funktion \\
Is there any difference in novice comprehension of a small program written in the event-driven and object-oriented styles? & 2002 & WoS, Scopus & \emph{citerad}: händelser från början \\
Learning computer science concepts with Scratch & 2013 & OA & \emph{citerad}: hur programmet styrs \\
Non-Event Driven Graphics API for Programming Education & 2013 & WoS & \emph{citerad}: händelsestyrd programmering som svårighet \\
Plan Composition Using Higher-Order Functions & 2023 & WoS, Scopus & \emph{citerad}: funktioner som värden \\
Students Struggle to Explain Their Own Program Code & 2021 & WoS & \emph{citerad}: lämna över eller anropa en funktion \\
Using HCI techniques to design a more usable programming system & 2002 & IEEE & \emph{citerad}: händelser från början \\
\enquote{Hear} and \enquote{Play} Students Misconceptions on Concurrent Programming using Sonic Pi & 2024 & OA & stöder påståendet (0.98) \\
Development and evaluation of a model of programming errors & 2003 & IEEE & stöder påståendet (0.92) \\
Extending the Object-Oriented Notional Machine Notation with Inheritance, Polymorphism, and GUI Events & 2017 & WoS, IEEE & stöder påståendet (0.96) \\
Isopleth & 2019 & OA & stöder påståendet (0.97) \\
Language and Temporal Aspects: A Qualitative Study on Trigger Interpretation in Trigger-Action Rules & 2023 & Scopus & stöder påståendet (0.98) \\
Making trigger-action rules more comprehensible: Investigating which linguistic clues effectively guide non-programmers & 2025 & Scopus & stöder påståendet (0.94) \\
Map, Filter, and Conquer: A Visual Tool for Learning Higher-Order Functions & 2025 & Scopus & stöder påståendet (0.93) \\
Six Learning Barriers in End-User Programming Systems & 2004 & OA & stöder påståendet (0.88) \\
Teaching Visual Programming to High School Students, Experience from Pakistan & 2005 & IEEE & stöder påståendet (0.97) \\
Tips for Teaching Types and Functions & 2008 & WoS & stöder påståendet (0.99) \\
Towards a Notional Machine for Runtime Stacks and Scope: When Stacks Don't Stack Up & 2023 & WoS, Scopus & stöder påståendet (0.96) \\
Understanding beginners' mistakes with Haskell & 2015 & OA & stöder påståendet (0.81) \\
Visual Tracing System for Novice Programmers of Java Applets Using Aspect-Oriented Programming & 2009 & IEEE & stöder påståendet (0.98) \\
A CS1 to CS2 bridge class using 2D game programming & 2006 & OA & kvalificerar eller motsäger påståendet (0.81) \\
A library to support a graphics-based object-first approach to CS 1 & 2001 & OA & kvalificerar eller motsäger påståendet (0.85) \\
Adding breadth to CS1 and CS2 courses through visual and interactive programming projects & 1999 & OA & kvalificerar eller motsäger påståendet (0.85) \\
Alice: Easy-to-Learn 3D Scripting for Novices & 1998 & OA & kvalificerar eller motsäger påståendet (0.82) \\
ConCodeIt! A Comparison of Concurrency Interfaces in Block-Based Visual Robot Programming & 2020 & IEEE & kvalificerar eller motsäger påståendet (0.98) \\
DoodlePad: next-gen event-driven programming for CS1 & 2017 & OA & kvalificerar eller motsäger påståendet (0.90) \\
Easy, realistic GUIs with Java in CS1 & 2000 & OA & kvalificerar eller motsäger påståendet (0.95) \\
Expressing Computer Science Concepts Through Kodu Game Lab & 2011 & WoS, Scopus & kvalificerar eller motsäger påståendet (0.96) \\
If when is better than if (and while might help): On the importance of influencing mental models in EUD (a pilot study) & 2020 & Scopus & kvalificerar eller motsäger påståendet (0.94) \\
Infusing an HtDP-based CS1 with distributed programming using functional video games & 2018 & OA & kvalificerar eller motsäger påståendet (0.86) \\
Introducing high school students to event driven programming & 1999 & IEEE & kvalificerar eller motsäger påståendet (0.89) \\
Introducing Java: the case for fundamentals-first & 2004 & WoS & kvalificerar eller motsäger påståendet (0.76) \\
Teaching Novices Programming and Important Applications in a Single Semester - Critical Factors from Zero to Portable GUI Programming in Four Hours & 2024 & Scopus & kvalificerar eller motsäger påståendet (0.86) \\
Using graphics to support the teaching of fundamental object-oriented principles in CS1 & 2003 & OA & kvalificerar eller motsäger påståendet (0.91) \\
Workshop \enquote{Business Programming} - Critical Factors from Zero to Portable GUI Programming in 4 Hours & 2023 & IEEE & kvalificerar eller motsäger påståendet (0.88) \\
\end{longtable}
\endgroup


\chapter{Är det en etablerad princip att hålla gränssnitt och
  programlogik isär, och gör det programmet lättare att pröva?}
\label{app:grafik-separation}
\chapterprecis{Författaren har ännu inte granskat resultaten i den här
  bilagan i sin helhet.}

% Author's note: the claim is made in sec:utan-fonster (subsection
% "Hälsningen utan fönster", contents.nw), wording fixed 2026-09-30:
% "Att hålla gränssnitt och logik isär är en etablerad designprincip,
% formulerad som Model--View--Controller\autocite{Reenskaug1979MVC,
% KrasnerPope1988} och ett fall av det Dijkstra kallade separation av
% ansvarsområden (separation of concerns)\autocite{Dijkstra1982Role}. Att
% logiken då går att pröva utan fönster är ett av de skäl som ges för
% principen, eftersom det är dyrt och osäkert att pröva ett program genom
% dess gränssnitt\autocite{Coppola2019Fragility,Hellmann2014AGT}."
% followed by the appendix-pointer footnote, which carries the
% qualification that "lättare att pröva och felsöka" is backed only
% indirectly.
I \cref{sec:utan-fonster} håller materialet funktionen
\mintinline{python}{greeting}, som bygger hälsningen, utanför klassen
som bygger fönstret, och kör den sedan utan fönster.  Påståendet som
följer är tvådelat.  Det första ledet säger att det är en
\emph{etablerad designprincip} att hålla ett programs gränssnitt och dess
logik isär, formulerad som Model--View--Controller, som Reenskaug
skrev ner vid Xerox PARC\autocite{Reenskaug1979MVC} och Krasner och
Pope publicerade för Smalltalk-80\autocite{KrasnerPope1988}, och ett fall
av det Dijkstra kallade separation av ansvarsområden
(\foreignlanguage{english}{separation of concerns})\autocite{%
Dijkstra1982Role}.  Det andra ledet säger vad principen är bra för: att
logiken då går att pröva utan fönster är ett av de skäl som ges för
principen, eftersom det är dyrt och osäkert att pröva ett program genom
dess gränssnitt.  Den starkare formen av nyttan, att isärhållandet gör
programmet lättare att testa och felsöka, påstår materialet inte i
texten; det är den formen kapitlet prövar, och det är den som fotnoten
vid påståendet sammanfattar med sina förbehåll.  Det första ledet rymmer
i sin tur två frågor, en om ursprung, som besvaras genom att läsa
källorna själva, och en om etablering, som kräver belägg för att
principen faktiskt tagits upp och lärs ut och en motsökning efter
kritik.  Det andra ledet är ett påstående om nytta, och det prövas
tvåsidigt: finns det belägg för att isärhållandet gör prövning och
felsökning lättare, och finns det belägg för motsatsen eller för att
nyttan är mindre än den sägs vara?  Frågan som ska besvaras är:
\emph{är det en etablerad princip att hålla gränssnitt och programlogik
isär, och gör det programmet lättare att pröva?}
\Cref{sec:sep-metod-grafik} beskriver hur påståendet prövats,
\cref{sec:sep-resultat-grafik} vad som fanns,
\cref{sec:sep-diskussion-grafik} vad som begränsar svaret,
\cref{sec:sep-slutsats-grafik} ger det och
\cref{sec:sep-fortsatt-grafik} säger vad som återstår att göra;
\cref{tab:sok-grafik-separation} sist i kapitlet listar de träffar som
rör påståendet.

Principen är släkt med två som redan är prövade i tidigare
föreläsningar: principen om ett enda ansvar, i föreläsningen
\emph{Funktioner}, bilaga~B, som fann den etablerad men fann att vad som
räknas som \emph{ett} ansvar är en bedömning, och inkapsling, i
föreläsningen \emph{Klasser och objekt}, bilaga~C, som fann den
etablerad sedan Parnas men löst definierad.  Det här kapitlet gör inte
om de sökningarna; det gäller den särskilda uppdelningen mellan
gränssnitt och logik och den särskilda nyttan att logiken kan prövas
för sig.

\section{Metod}
\label{sec:sep-metod-grafik}

Ursprunget söktes som kända dokument, eftersom inget av dem finns i
databaserna: Reenskaugs anteckning från december 1979 är ett
maskinskrivet internt dokument, som lästes i den kopia han själv lagt
ut; Krasner och Popes artikel kom ut 1988 i en tidskrift utan
\textsc{doi}-nummer och lästes i en kopia av tidskriftssidorna; och
Dijkstras essä från 1974 lästes i den utskrift som hans arkiv vid
University of Texas tillhandahåller, medan den tryckta versionens
uppgifter kontrollerades mot förlagets register.  Läroboken
\emph{Design Patterns}\autocite{Gamma1994DesignPatterns} lästes i
författarens eget exemplar, och där, bland författarens egna böcker,
fanns också Holubs kritik av mönstret\autocite{Holub2004}, som ingen
av sökningarna gav.  Några av de källor som bär
praktikerperspektivet, Fowlers två texter om gränssnittsarkitektur,
Feathers artikel om den \enquote{ödmjuka dialogrutan} och Meszaros
mönster \foreignlanguage{english}{Humble Object}, är webbtexter som
ingen av databaserna indexerar; de lästes där de publicerats och
behandlas som det de är, erfarenhet och inte mätningar.

Etableringen och nyttan söktes ämnesdrivet och tvåsidigt den 30
september 2026 med verktyget \texttt{scholar}, mot de databaser som
svarade: OpenAlex, IEEE Xplore, Web of Science och Scopus.  DBLP
svarade den dagen på varje fråga antingen med en sida som kontrollerar
att besökaren inte är en robot eller med en nerkopplad förbindelse, och
Semantic Scholar avvisade sin nyckel; båda saknas.  Stödfrågorna
S1--S9 söker Model--View--Controller och dess efterföljare
(Model--View--Presenter, Model--View--ViewModel, \foreignlanguage{english}{%
Presentation Model}, \foreignlanguage{english}{Humble Object}),
separation av ansvarsområden mellan gränssnitt och affärslogik,
dialogoberoende (\foreignlanguage{english}{dialogue independence}), som
är människa--dator-forskningens namn på samma idé, prövbarhet och
felsökning, kostnaden för att pröva genom gränssnittet, undervisning,
och jämförelser av mönstren i appar för Android och iOS.  Motfrågorna
M1--M8 använder samma begrepp tillsammans med ord för kritik,
nackdelar, overhead, missförstånd, kodlukter och brott mot mönstret,
uteblivna skillnader, fel som sitter i gränssnittet själv och test av
logiken som inte räcker.  \Cref{tab:sok-queries-grafik-separation} visar
frågorna och antalet träffar per databas.

% Author's note (verification and reproduction): see the block above the
% hit-list \input at the end of this chapter.  The query table is
% transcribed from matrix.tsv / queries.json in the research directory
% (scholar session grafik-separation), which holds the exact string sent
% to each provider.
\begin{table}[htbp]
  \centering
  \caption{Sökfrågor och antal träffar per databas, isärhållandet av
    gränssnitt och logik, 30 september 2026, verktyget
    \texttt{scholar}.  S = stödfråga, M = motfråga.  Databaser: OpenAlex
    (OA; inte \foreignlanguage{english}{open access}), IEEE Xplore
    (IEEE), Web of Science (WoS), Scopus.  Högst 40 träffar hämtades per
    databas och fråga, så 40 betyder \enquote{minst 40}.  Frågan visas
    i förkortad form; den skickades som boolesk fritext till OpenAlex och
    IEEE, som \texttt{TS=(\dots)} till Web of Science och som
    \texttt{TITLE-ABS-KEY(\dots)} till Scopus, med samma nyckelord
    överallt.  MVC, MVP och MVVM står för mönstren Model--View--Controller,
    Model--View--Presenter och Model--View--ViewModel, UIMS för
    \foreignlanguage{english}{user interface management system}.  DBLP
    svarade inte och Semantic Scholar avvisade sin nyckel; de saknas.}
  \label{tab:sok-queries-grafik-separation}
  \footnotesize
  \begin{tabular}{@{}lp{0.62\linewidth}rrrr@{}}
    \toprule
    & Nyckelord & OA & IEEE & WoS & Scopus \\
    \midrule
    S1 & \enquote{model-view-controller} \textsc{and} (\enquote{user
         interface} \textsc{or} GUI) \textsc{and} (separation
         \textsc{or} decoupling)
       & 40 & 18 & 21 & 26 \\
    S2 & \enquote{separation of concerns} \textsc{and} (\enquote{user
         interface} \textsc{or} GUI \textsc{or} presentation)
         \textsc{and} (\enquote{business logic} \textsc{or}
         \enquote{application logic} \textsc{or} \enquote{domain logic}
         \textsc{or} model)
       & 39 & 29 & 40 & 40 \\
    S3 & (MVP \textsc{or} MVVM \textsc{or} \enquote{presentation model}
         \textsc{or} \enquote{humble object}) \textsc{and} (\enquote{user
         interface} \textsc{or} GUI) \textsc{and} (architecture
         \textsc{or} pattern)
       & 40 & 26 & 28 & 40 \\
    S4 & (\enquote{user interface} \textsc{or} GUI \textsc{or}
         \enquote{presentation layer}) \textsc{and} (\enquote{business
         logic} \textsc{or} \enquote{application logic} \textsc{or}
         \enquote{domain logic}) \textsc{and} (separation \textsc{or}
         decoupling) \textsc{and} (testability \textsc{or} testing
         \textsc{or} \enquote{unit test} \textsc{or} debugging)
       & 40 & 23 & 9 & 6 \\
    S5 & \enquote{GUI testing} \textsc{and} (difficult \textsc{or}
         challenges \textsc{or} cost \textsc{or} expensive \textsc{or}
         fragile \textsc{or} flaky \textsc{or} maintenance)
       & 40 & 40 & 40 & 40 \\
    S6 & (MVC \textsc{or} MVP \textsc{or} MVVM \textsc{or}
         \enquote{separation of concerns}) \textsc{and} (\enquote{user
         interface} \textsc{or} GUI \textsc{or} presentation)
         \textsc{and} (maintainability \textsc{or} testability
         \textsc{or} \enquote{empirical study} \textsc{or} experiment
         \textsc{or} evaluation)
       & 40 & 40 & 40 & 40 \\
    S7 & (MVC \textsc{or} \enquote{separation of concerns})
         \textsc{and} (\enquote{user interface} \textsc{or} GUI)
         \textsc{and} (teaching \textsc{or} education \textsc{or}
         students \textsc{or} novices \textsc{or} CS1)
       & 40 & 29 & 23 & 40 \\
    S8 & (UIMS \textsc{or} \enquote{Seeheim model} \textsc{or}
         \enquote{presentation-abstraction-control} \textsc{or}
         \enquote{dialogue independence}) \textsc{and} (separation
         \textsc{or} separability \textsc{or} independence \textsc{or}
         decoupling)
       & 40 & 27 & 25 & 40 \\
    S9 & (MVC \textsc{or} MVP \textsc{or} MVVM \textsc{or} \enquote{clean
         architecture}) \textsc{and} (Android \textsc{or} mobile
         \textsc{or} apps) \textsc{and} (testability \textsc{or} testing
         \textsc{or} \enquote{unit test} \textsc{or} maintainability)
         \textsc{and} (empirical \textsc{or} study \textsc{or}
         developers \textsc{or} practitioners)
       & 40 & 40 & 34 & 40 \\
    \midrule
    M1 & (MVC \textsc{or} MVP \textsc{or} MVVM \textsc{or}
         \enquote{separation of concerns}) \textsc{and} (\enquote{user
         interface} \textsc{or} GUI) \textsc{and} (criticism \textsc{or}
         drawbacks \textsc{or} limitations \textsc{or} misunderstood
         \textsc{or} \enquote{anti-pattern} \textsc{or} overhead
         \textsc{or} \enquote{over-engineering} \textsc{or} complexity)
       & 40 & 23 & 27 & 40 \\
    M2 & (MVC \textsc{or} MVP \textsc{or} MVVM \textsc{or}
         \enquote{separation of concerns}) \textsc{and} (testability
         \textsc{or} testing \textsc{or} maintainability \textsc{or}
         debugging) \textsc{and} (\enquote{no significant} \textsc{or}
         \enquote{no difference} \textsc{or} \enquote{does not}
         \textsc{or} worse \textsc{or} violations \textsc{or} smells)
       & 40 & 40 & 40 & 40 \\
    M3 & som S4 men med (cost \textsc{or} overhead \textsc{or}
         complexity \textsc{or} drawbacks \textsc{or} limitations
         \textsc{or} \enquote{not sufficient} \textsc{or} tangled
         \textsc{or} mixing \textsc{or} violations) i stället för
         prövbarhetsorden
       & 40 & 31 & 11 & 13 \\
    M4 & (\enquote{GUI faults} \textsc{or} \enquote{GUI bugs} \textsc{or}
         \enquote{GUI defects} \textsc{or} \enquote{user interface
         faults}) \textsc{and} (empirical \textsc{or} study \textsc{or}
         taxonomy \textsc{or} classification)
       & 38 & 7 & 5 & 16 \\
    M5 & \enquote{GUI testing} \textsc{and} (\enquote{business logic}
         \textsc{or} \enquote{application logic} \textsc{or} \enquote{unit
         testing}) \textsc{and} (insufficient \textsc{or} complement
         \textsc{or} \enquote{end-to-end} \textsc{or} \enquote{system
         testing} \textsc{or} missed \textsc{or} limitations)
       & 38 & 12 & 5 & 11 \\
    M6 & \enquote{separation of concerns} \textsc{and} (criticism
         \textsc{or} limitations \textsc{or} \enquote{crosscutting
         concerns} \textsc{or} tangling \textsc{or} scattering)
         \textsc{and} (maintainability \textsc{or} testability
         \textsc{or} \enquote{empirical study} \textsc{or} experiment)
       & 40 & 10 & 22 & 36 \\
    M7 & som S8 och därtill (criticism \textsc{or} limitations
         \textsc{or} drawbacks \textsc{or} problems \textsc{or} failure
         \textsc{or} \enquote{semantic feedback})
       & 40 & 6 & 5 & 13 \\
    M8 & som S9 men med (\enquote{no significant} \textsc{or} \enquote{no
         difference} \textsc{or} worse \textsc{or} overhead \textsc{or}
         drawbacks \textsc{or} \enquote{trade-offs} \textsc{or}
         violations \textsc{or} smells) i stället för studieorden
       & 40 & 4 & 3 & 7 \\
    \bottomrule
  \end{tabular}
\end{table}

Sammanlagt lämnade databaserna 1\,946 poster.  Flera databaser lämnar
samma arbete, och två poster räknas som ett arbete när titlarna är
desamma bortsett från versaler och skiljetecken; efter den
sammanslagningen återstår 1\,185 arbeten, och det är arbeten som räknas
här och i \cref{tab:sok-grafik-separation}.  Alla klassades: först med
en språkmodell mot en beskrivning av frågan och fyra kategorier ---
stöder påståendet, kvalificerar eller motsäger det, angränsande
delämne, annat ämne --- och sedan för hand: varje post som modellen
helt eller delvis satt i den andra kategorin, och varje post med
konfidens under 0,7, sammanlagt 88 poster, lästes i sammanfattning, och
22 av dem flyttades till angränsande delämne, oftast verktyg för
gränssnittstest och studier av prestanda som inte säger något om
isärhållandet.  Efter granskningen citeras 28 poster, 205 stöder
påståendet och 25 kvalificerar eller motsäger det; resten är
angränsande (631) eller felträffar (296).  Antalet stödjande är en
övre gräns, eftersom modellen gärna räknar ett program som bara
använder Model--View--Controller som stöd, och de raderna lästes inte
alla för hand.  Av de 36 källor kapitlet citerar lästes 14 i fulltext;
för de övriga står det i källfilen att bara sammanfattningen lästs, och
\cref{sec:sep-diskussion-grafik} återkommer till det.  Hur varje källa
hittades och verifierades är antecknat på det sätt föreläsningen
\emph{Algoritmiskt tänkande}, bilaga~A, beskriver.

\section{Resultat}
\label{sec:sep-resultat-grafik}

\paragraph{Ursprunget.}  Reenskaugs anteckning delar ett program i
tre roller.  Modellen representerar kunskap, och vyn är dess bild:
\textcquote[1]{Reenskaug1979MVC}{A view is a (visual) representation of
its model. [\dots] It is thus acting as a presentation filter.}
Kontrollern tar hand om användarens inmatning, och vyn ska hållas fri
från den: \textcquote[1]{Reenskaug1979MVC}{a view should never know
about user input, such as mouse operations and keystrokes}.
\Textcite{KrasnerPope1988} beskrev nio år senare samma uppdelning som
Smalltalk-80:s sätt att bygga interaktiva program, och motiverade den
med modularitet: \textcquote[26]{KrasnerPope1988}{Isolating functional
units from each other as much as possible makes it easier for the
application designer to understand and modify each particular unit,
without having to know everything about the other units.}  Den
tredelning de beskriver innebär att skilja
\textcquote[26]{KrasnerPope1988}{(1) the parts that represent the model
of the underlying application domain from (2) the way the model is
presented to the user and from (3) the way the user interacts with it}.
Värt att notera är vad ursprungskällorna \emph{inte} säger: de nyttor
de räknar upp är att förstå, ändra och återanvända delarna, inte att
pröva dem.  Uttrycket \enquote{separation of concerns} är Dijkstras, men
hos honom gäller det ett sätt att tänka, att studera en aspekt i taget
--- korrekthet en dag, effektivitet en annan --- och inte just
gränssnitt och logik: \textcquote{Dijkstra1982Role}{It is what I
sometimes have called \enquote{the separation of concerns}, which, even
if not perfectly possible, is yet the only available technique for
effective ordering of one's thoughts, that I know of.}  Att dela ett
program mellan gränssnitt och logik är en senare tillämpning av den
tanken.

\paragraph{Etableringen.}  Principen fick fäste i tre traditioner.  I
mönsterlitteraturen inleder \textcite{Gamma1994DesignPatterns} sin bok
med Model--View--Controller som exempel på vad ett mönster är:
\textcquote[4]{Gamma1994DesignPatterns}{Before MVC, user interface
designs tended to lump these objects together. MVC decouples them to
increase flexibility and reuse.}  Mönstret fick efterföljare:
\textcite{Potel1996MVP} beskriver Model--View--Presenter som en
generalisering av det, vars första steg var att
\textcquote[2]{Potel1996MVP}{formalize the separation between the Model
and the View-Controller}, och konstaterar att Model--View--Controller
redan \textcquote[1]{Potel1996MVP}{has been reused and adopted to
varying degrees in most other GUI class libraries and application
frameworks}; \textcite{SyromiatnikovWeyns2014} kartlägger hela familjen
och kallar Model--View--Controller \textcquote{SyromiatnikovWeyns2014}{a
common design pattern to integrate a user interface with the
application domain logic}.  I människa--dator-forskningen hette samma
idé dialogoberoende, och i Hartson och Hix
översikt\autocite{HartsonHix1989} är den
\textcquote{HartsonHix1989}{the keystone concept upon which all the
other concepts depend}, en egenskap som
\textcquote{HartsonHix1989}{separates design of the interface from
design of the computational component of an application system so that
modifications in either tend not to cause changes in the other}; samma
tid framhölls den som en väg till mer lätthanterlig programvara
\parencite{YipRobson1991}, och inom programvaruarkitekturen var
användbarhet länge inte arkitektens sak
\textcquote{John2005Usability}{beyond separating the user interface
from the remainder of the application}.  Bland dagens praktiker beskriver
\textcite{Fowler2006GUIArchitectures} det han kallar
\foreignlanguage{english}{Separated Presentation} som
\textcquote{Fowler2006GUIArchitectures}{the idea that was the most
influential to later frameworks}, och menar att stilen
\textcquote{Fowler2006SeparatedPresentation}{came into vogue with
Model-View-Controller architecture and is widely used}; i en studie av
hur appar för Android byggs var Model--View--Presenter det mönster
som studiens källor oftast tog upp, och riktlinjerna säger att affärslogiken
ska hållas helt frikopplad från Android-ramverket
\parencite{Verdecchia2019Android}.  Mönstret lärs också ut:
\textcite{Tao2002MVC} beskriver hur Model--View--Controller undervisas
som arkitekturen för grafiska program, och \textcite{HansenFossum2005}
att det ofta lärs ut som föreskrift, med tre sorters objekt att bygga
grafiska program av.  Separation av ansvarsområden i allmänhet
tillskrivs, \textcquote{Tarr1999NDegrees}{[d]one well}, lägre
komplexitet, bättre återanvändning och enklare vidareutveckling
\parencite{Tarr1999NDegrees}.

\paragraph{Nyttan: att pröva logiken utan gränssnittet.}  Att
prövbarheten är ett uttalat skäl för principen syns tydligast i
praktikerlitteraturen.  \Textcite{Feathers2002Humble} beskriver hur
logik i en dialogruta leder till att
\textcquote[1]{Feathers2002Humble}{you introduce a bug because you
couldn't get your dialog class into a test harness}, och hur den som
flyttar logiken till ett eget objekt får kvar
\textcquote[8]{Feathers2002Humble}{a very thin wedge at the interface
that consists of simple delegation and value setting} som enda oprövade
kod.  \Textcite{Meszaros2007Humble} gjorde av det ett allmänt mönster
för kod som är svår att pröva för att den sitter fast i sin omgivning:
\textcquote{Meszaros2007Humble}{We extract the logic into a separate
easy-to-test component that is decoupled from its environment.}
\Textcite{Fowler2006GUIArchitectures} sammanfattar rörelsen: genom att
pröva presentatören eller presentationsmodellen
\textcquote{Fowler2006GUIArchitectures}{you test most of the risk of the
UI without having to touch the hard-to-test widgets}, och ett
industriellt team rapporterar att sådana test är
\textcquote{Alles2006PresenterFirst}{an economical alternative to
automated GUI system tests}.  Bland Android-praktikernas riktlinjer är
skälet detsamma: frikopplade komponenter ger
\textcquote[6]{Verdecchia2019Android}{higher testability of the core
logic of the app by making components unit-testable (ideally without
requiring an emulator)}.

Det som har mätts är den andra sidan av saken, att det är dyrt och
osäkert att pröva genom gränssnittet.  Att pröva grafiska gränssnitt för
hand beskrivs som \textcquote{Hellmann2014AGT}{expensive and complex},
och i samma intervjustudie berättar en deltagare att ett misslyckat
gränssnittstest \textcquote[3]{Hellmann2014AGT}{could not show us what
is the root of the problem, whether there is some problem in the
business logic of the application or something wrong with the user
interface} --- det är felsökningsledet i påståendet, sett från andra
hållet.  Gränssnittstest är ömtåliga: i Android-projekt med öppen
källkod var i genomsnitt 70~procent av ändringarna i dem orsakade av
ändringar i gränssnittet och inte av fel \parencite{Coppola2019Fragility}.
Och de når inte all logik: röktest genom gränssnittet kan inte täcka
vissa delar av programkoden \parencite{MemonXie2005}, och i ett
experiment prövade gränssnittstest programlogiken mindre än test som
anropar logiken direkt \parencite{AlkhalidLabiche2017}; att bara utforska
gränssnittet missar ofta mycket av programmets beteende
\parencite{Gross2012EXSYST}.  Direkta
jämförelser av mönstrens prövbarhet finns, men de är få och små.
\Textcite{Lou2016Android} konstaterar att fördelen hos
Model--View--Presenter och Model--View--ViewModel
\textcquote{Lou2016Android}{is claimed [\dots] But there is no empirical
data to support this point of view}, och finner den sedan i en analys av
en enda app; \textcite{MagicsVerkman2023iOS} skriver likaså att
\textcquote{MagicsVerkman2023iOS}{there are few empirical data that
support a specific choice} och jämför tre mönster i en fallstudie med
en app; \textcite{Wilson2022MVVM} jämför prövbarheten med mätetal på en
exempelkodbas, och en fjärde studie fann att en app byggd med
Model--View--Presenter var dubbelt så modulär som en utan mönster, utan
att mäta prövbarheten\autocite{Prabowo2018MVP}.

\paragraph{Motsökningen.}  Ingen studie som mätt nyttan och funnit den
utebli kom fram.  Det som kom fram kvalificerar påståendet på fyra
sätt.  För det första prövar ett test av logiken inte gränssnittet.  I
två industriella system låg 60~procent av felen som kunderna
rapporterade i gränssnittet och inte i själva tillämpningen, och de
tog längre tid att rätta \parencite{RobinsonBrooks2009}; gränssnittskod
har egna slags fel, som kräver egna test \parencite{Lelli2015GUIDefects},
och test som genereras direkt från koden kan bli orealistiska eller
irrelevanta \parencite{Gross2012EXSYST}.
Isärhållandet flyttar alltså logiken dit den är lätt att pröva, men det
tar inte bort behovet av att pröva gränssnittet, och det räcker inte
heller för allt ett gränssnitt ska klara: flera användbarhetskrav
kräver mer av arkitekturen än isärhållandet ensamt ger
\parencite{John2005Usability}.  För det andra kostar
isärhållandet något.  Att hålla de två delarna i takt är ett problem i
sig \parencite{SotirovskiKruchten1995}, i statiskt typade språk kräver
mönstret avvägningar som \textcquote{VeitHerrmann2003}{either blur the
clear structure, destroy the intended independence, or introduce undue
administrative overhead}, och en genomgång av Model--View--ViewModel
räknar upp 15 nackdelar vid sidan av fördelarna och konstaterar att
mönstret ofta används fel i brist på vägledning
\parencite{Fuksa2025MVVM}.  Den tidiga människa--dator-forskningens
system byggda på isärhållande blev heller inte den framgång man väntat
sig: \textcquote{RathnamMannino1991}{Their adoption has been slow and at
times painful}, och ett rent isärhållande var ett mål som var svårt att
nå, eftersom gränssnitt där användaren direkt manipulerar objekt
kräver ständig återkoppling från tillämpningen \parencite{Hill1992ALV}.
För det tredje hålls isärhållandet inte alltid i praktiken.  I
Model--View--Controller-program för webben fann
\textcite{Aniche2018MVCSmells} att logik hamnar i fel lager --- en
utvecklare förklarade att det egna teamet
\textcquote[2128]{Aniche2018MVCSmells}{does not unit test CONTROLLERS,
and thus, complex logic and control flow in them should be avoided} ---
och att klasser med sådana brister ändrades oftare; men när klassernas
storlek hölls konstant fanns \textcquote[2139]{Aniche2018MVCSmells}{no
significant difference} i hur ofta de rättades för fel.  För det fjärde
är mycket av det som sägs om mönstren omdömen snarare än mätningar.
Model--View--Controller är enligt
\textcite{Fowler2006GUIArchitectures}
\textcquote{Fowler2006GUIArchitectures}{one
of the most misunderstood architectural patterns around}, och samme
författare skriver att svårigheten att pröva gränssnitt
\textcquote{Fowler2006GUIArchitectures}{is over-stated}
ibland.  \Textcite{Holub2004} går längre och menar att
Model--View--Controller i sin vanliga form
\textcquote[15]{Holub2004}{fails miserably as an application-level
architecture because MVC requires the controller to know way too much
about how the model-level objects are implemented} --- men hans
invändning gäller hur kontrollern kopplas till modellen, inte att
presentation och logik hålls isär, och han vill ha mer inkapsling, inte
mindre.  Om arkitektur för Android-appar säger
\textcite{Verdecchia2019Android} att den
\textcquote[1]{Verdecchia2019Android}{is
still debated and subject to conflicting opinions usually influenced by
technological hypes rather than objective evidence}.  Och i allmän form
har separation av ansvarsområden en känd gräns: de vanliga metoderna
delar ett program efter en enda dimension, och de ansvarsområden som
skär tvärs över den blir svåra att hålla isär \parencite{Tarr1999NDegrees}.

\section{Diskussion och begränsningar}
\label{sec:sep-diskussion-grafik}

Fyra saker begränsar svaret.  Belägget för nyttan är av två slag, och
inget av dem är ett kontrollerat försök.  Det ena är praktikers
erfarenhet --- Feathers, Meszaros, Fowler, teamet bakom
\foreignlanguage{english}{Presenter First} och Android-praktikerna ---
som är samstämmig men inte mätt.  Det andra är mätningar av att prövning
genom gränssnittet är dyr, ömtålig och når mindre av logiken, vilket gör
det rimligt att det lönar sig att kunna pröva logiken för sig, men som
inte mäter just det.  De få studier som jämför mönstrens prövbarhet
bygger på en eller två appar och på mätetal som ersätter själva
prövningen.  Källorna gäller dessutom Java, C\#, Smalltalk, Android och
iOS, yrkesverksamma utvecklare och stora program; ingen gäller Python,
\mintinline{python}{tkinter} eller nybörjare, och det program materialet
visar är litet.  Att logiken går att köra utan fönster är ändå inte ett
empiriskt påstående: det följer av hur programmet är byggt, och
\cref{sec:utan-fonster} visar det i en körning.  Verifieringen är
ojämn: Reenskaug, Krasner och Pope, Dijkstras essä, \emph{Design
Patterns}, Holub, Potel, Fowler, Feathers, Hellmann med flera, Verdecchia
med flera, Aniche med flera och Lelli med flera lästes i fulltext ---
de två böckerna i författarens egna exemplar --- och Meszaros mönster
i den förhandsversion på webben som föregick boken, medan de övriga
bara kunde läsas i sammanfattning, eftersom förlagen avvisar
automatisk hämtning.  Sidnumren i Aniche med fleras artikel är räknade
fram från den öppna kopian, som saknar sidnummer.  Slutligen är
täckningen ofullständig: DBLP och Semantic Scholar saknas, träffantalen i
\cref{tab:sok-queries-grafik-separation} är tak, praktikerlitteraturen
finns inte i databaserna och hämtades som kända dokument, och
klassningen är modellstödd, så antalet stödjande poster är en övre gräns.

\section{Slutsats}
\label{sec:sep-slutsats-grafik}

Ja på den första frågan, och på den andra med förbehåll.  Att hålla
gränssnitt och programlogik isär är en etablerad designprincip: den
formulerades i Model--View--Controller
\parencite{Reenskaug1979MVC,KrasnerPope1988}, bär namn och status i
mönsterlitteraturen \parencite{Gamma1994DesignPatterns,Potel1996MVP},
i människa--dator-forskningen \parencite{HartsonHix1989} och bland
dagens praktiker \parencite{Fowler2006SeparatedPresentation,
Verdecchia2019Android}, och den lärs ut \parencite{Tao2002MVC,
HansenFossum2005}.  Uttrycket separation av ansvarsområden är
Dijkstras \parencite{Dijkstra1982Role}, men hos honom gäller det
tänkandet i allmänhet; uppdelningen mellan gränssnitt och logik är en
tillämpning av det.  Att logiken då kan köras och prövas utan
gränssnittet följer av konstruktionen, och det är det skäl
praktikerlitteraturen ger för principen
\parencite{Feathers2002Humble,Meszaros2007Humble,%
Fowler2006GUIArchitectures}.
Att det gör programmet lättare att pröva och felsöka stöds indirekt:
prövning genom gränssnittet är dyr och ömtålig, når mindre av logiken och
säger inte var felet sitter
\parencite{Hellmann2014AGT,Coppola2019Fragility,AlkhalidLabiche2017}.
Förbehållen är att ingen kontrollerad studie har mätt själva vinsten ---
de jämförelser som finns gäller en app i taget
\parencite{Lou2016Android,MagicsVerkman2023iOS} --- att ursprungskällorna
motiverar principen med förståelse, ändring och återanvändning snarare
än med prövning, och att isärhållandet inte gör det onödigt att pröva
gränssnittet, där en stor del av felen sitter
\parencite{RobinsonBrooks2009}.  Materialet kan därför säga att
principen är etablerad, att logiken går att pröva utan fönster och att
det är ett av de skäl som ges för principen, eftersom det är dyrt och
osäkert att pröva genom gränssnittet; det bör inte säga att programmet
mätbart blir lättare att felsöka.

\section{Fortsatt arbete}
\label{sec:sep-fortsatt-grafik}

Sökningen lämnade tre saker ogjorda.  Två databaser saknas helt: DBLP
svarade inte och Semantic Scholar avvisade sin nyckel
(\cref{sec:sep-metod-grafik}).  Samma 17 frågor bör ställas där när
tjänsterna svarar igen; DBLP indexerar just den konferenslitteratur om
programvaruarkitektur och testning som frågorna söker.  Flera källor är
därtill lästa bara i sammanfattning.  Tyngst väger de fyra som bär
nyttoledet\autocite{AlkhalidLabiche2017,Coppola2019Fragility,%
MemonXie2005,Lou2016Android}: en läsning i fulltext skulle visa hur stora
skillnaderna var och i vilka program de uppmättes.  Wilson med fleras
jämförelse\autocite{Wilson2022MVVM} säger i sammanfattningen inte vad den
fann, och Robinson och Brooks siffra om var felen sitter\autocite{%
RobinsonBrooks2009} gäller två industriella system och bör läsas i sitt
sammanhang.  Två kritiska arbeten om isärhållandet i den tidiga
människa--dator-forskningen, av Cockton från 1987\footnote{Gilbert
Cockton, \enquote{Interaction ergonomics, control and separation: open
problems in user interface management}, \emph{Information and Software
Technology} 29(4), 1987, s.~176--191, \textsc{doi}
10.1016/0950-5849(87)90066-8.} och
av Reynolds från 1997\footnote{Carson Reynolds, \enquote{A critical
examination of separable user interface management systems}, \emph{ACM
SIGCHI Bulletin} 29(3), 1997, s.~41--45, \textsc{doi}
10.1145/264853.264874.}, kom fram i motsökningen men kunde inte läsas,
och Meszaros bok bör kontrolleras mot den webbversion som citeras här.

För ett säkrare svar saknas ett slags studie.  Ingen av källorna har
mätt det påståendet faktiskt säger: att ett program med logiken skild
från gränssnittet är lättare att pröva och felsöka än samma program utan
den uppdelningen.  Ett kontrollerat försök där studenter får hitta och
rätta ett inplanterat fel i två versioner av samma lilla
\mintinline{python}{tkinter}-program, den ena med logiken i en egen
funktion och den andra med logiken i fönsterklassen, skulle mäta just
det, i den population materialet vänder sig till.  Faller ett sådant
försök ut till uppdelningens fördel kan materialet säga att programmet
blir lättare att felsöka, inte bara att logiken går att pröva för sig;
faller det ut åt andra hållet bör materialet motivera uppdelningen med
återanvändning och läsbarhet, som ursprungskällorna gör, snarare än med
prövning.

\section{Träffar som rör påståendet}
\label{sec:sep-traffar-grafik}

\Cref{tab:sok-grafik-separation} listar de poster som rör påståendet
--- citerade, stödjande och kvalificerande --- med skälet till varje
beslut; de angränsande och de felträffar som uteslöts anges som antal i
tabelltexten, och antalet träffar per fråga står i
\cref{tab:sok-queries-grafik-separation}.
% Author's note (verification and reproduction):
%   scholar session : grafik-separation
%   exports         : litteratursokning/grafik-separation.csv,
%                     litteratursokning/grafik-separation-full.tex
%   searched        : 2026-09-30; providers openalex ieee wos scopus
%                     (dblp: every call reset or answered with an Anubis
%                     bot-check page, on dblp.org and both mirrors --
%                     absent; s2: key rejected HTTP 403 -- absent);
%                     limit 40 per provider and query; openalex M6 first
%                     hit HTTP 429 and was rerun; exact strings per
%                     provider in queries.json / matrix.tsv / raw/ of the
%                     research directory and in the session store
%   classification  : scholar llm context + scholar llm classify
%                     --no-examples --no-enrich (batches of 20), model
%                     openai-codex/gpt-5.6-sol; human review of all 88
%                     rows with a qualifies-claim tag or confidence < 0.7
%                     (abstracts); 22 moved to adjacent-subtopic by
%                     scholar sessions decide --discard -t adjacent-subtopic
%   theme tags      : scholar sessions decide --keep -t <category>
%                     -t <theme>: etablering, provbarhet, gui-test,
%                     gui-fel, kostnad
%   hit-list table  : scholar sessions export grafik-separation -f table
%                     --bearing-only --lang sv
%                     --label sok-grafik-separation
%                     --track "isärhållandet av gränssnitt och programlogik"
%                     --theme etablering="principen etablerad"
%                     --theme provbarhet="mönstrens prövbarhet"
%                     --theme gui-test="att pröva genom gränssnittet"
%                     --theme gui-fel="fel i själva gränssnittet"
%                     --theme kostnad="kostnader, brister och kritik"
%                     -o litteratursokning/grafik-separation-full
%   provenance      : ltnotes.bib (from grafik-separation.bib), 36 keys,
%                     each with CLAIM / FOUND-VIA / PICKED / QUOTE /
%                     VERIFIED / COUNTER / DATE; Gamma1994DesignPatterns
%                     already in ltnotes.bib: merge the block
%   known items     : Reenskaug1979MVC, KrasnerPope1988, Dijkstra1982Role,
%                     Gamma1994DesignPatterns (reMarkable), Holub2004
%                     (reMarkable, found by name search there, not by the
%                     topic search), Fowler2006GUIArchitectures,
%                     Fowler2006SeparatedPresentation, Feathers2002Humble,
%                     Meszaros2007Humble (web draft) -- not in the table
%   full text       : the nine known items above, Potel1996MVP,
%                     Hellmann2014AGT, Verdecchia2019Android (author copy),
%                     Aniche2018MVCSmells (OA; pages inferred 2120+n),
%                     Lelli2015GUIDefects
%   abstract-level  : the other 22 keys -- for the user: Alkhalid &
%                     Labiche 2017, Coppola 2019, Memon & Xie 2005,
%                     Lou 2016, Wilson 2022, Robinson & Brooks 2009 first
% --- scholar LaTeX fragment --------------------------------------------
% Generated by: scholar sessions export -f table
% Session: grafik-separation
% \input this file from the body of a document whose
% preamble loads:
%   \usepackage{longtable}
%   \usepackage{booktabs}
%   \usepackage{array}
% ----------------------------------------------------------------------

\begingroup\footnotesize
\begin{longtable}{@{}%
  >{\raggedright\arraybackslash}p{0.44\textwidth}c%
  >{\raggedright\arraybackslash}p{0.13\textwidth}%
  >{\raggedright\arraybackslash}p{0.26\textwidth}@{}}
\caption{Träffar som rör påståendet, isärhållandet av gränssnitt och programlogik (28 citerade, 205 som stöder påståendet och 25 som kvalificerar eller motsäger det, av 1185 unika träffar; de 631 angränsande och 296 felträffarna listas inte). Skälet till varje inkludering står i sista kolumnen; för maskinklassade rader anges modellens konfidens. Databaser: OpenAlex (OA; inte open access), IEEE Xplore (IEEE), Web of Science (WoS), Scopus.}\label{tab:sok-grafik-separation}\\
\toprule
Titel & År & Leverantör & Skäl \\
\midrule
\endfirsthead
\multicolumn{4}{@{}l}{\emph{\tablename~\thetable{} (forts.)}}\\
\toprule
Titel & År & Leverantör & Skäl \\
\midrule
\endhead
\midrule \multicolumn{4}{r@{}}{\emph{forts.\ på nästa sida}}\\
\endfoot
\bottomrule
\endlastfoot
A Comparison of Android Native App Architecture -- MVC, MVP and MVVM & 2016 & OA & \emph{citerad}: mönstrens prövbarhet \\
A Comparison of Architectural Patterns for Testability and Performance Quality for iOS Mobile Applications Development & 2023 & Scopus & \emph{citerad}: mönstrens prövbarhet \\
A Journey through the Land of Model-View-Design Patterns & 2014 & IEEE & \emph{citerad}: principen etablerad \\
A Tale of Two Development Approach: Empirical Study on The Maintainability and Modularity of Android Mobile Application with Anti-Pattern and Model-View-Presenter Design Pattern & 2018 & WoS, OA & \emph{citerad}: mönstrens prövbarhet \\
An Empirical Evaluation and Comparison of the Impact of MVVM and MVC GUI Driven Application Architectures on Maintainability and Testability & 2022 & IEEE & \emph{citerad}: mönstrens prövbarhet \\
An Exploratory Study of Automated GUI Testing: Goals, Issues, and Best Practices & 2014 & OA & \emph{citerad}: att pröva genom gränssnittet \\
An initial study of customer-reported GUI defects & 2009 & Scopus & \emph{citerad}: fel i själva gränssnittet \\
Bringing usability concerns to the design of software architecture & 2005 & WoS, Scopus & \emph{citerad}: principen etablerad \\
Classifying and Qualifying GUI Defects & 2015 & WoS, Scopus & \emph{citerad}: fel i själva gränssnittet \\
Code smells for Model-View-Controller architectures & 2017 & OA & \emph{citerad}: kostnader, brister och kritik \\
Component- vs. application-level MVC architecture & 2002 & IEEE & \emph{citerad}: principen etablerad \\
EXSYST: Search-based GUI testing & 2012 & IEEE & \emph{citerad}: att pröva genom gränssnittet \\
Guidelines for Architecting Android Apps: A Mixed-Method Empirical Study & 2019 & OA & \emph{citerad}: principen etablerad \\
How Does GUI Testing Exercise Application Logic Functionality? & 2017 & Scopus & \emph{citerad}: att pröva genom gränssnittet \\
Human-Computer Interface Development: Concepts and Systems for Its Management & 1989 & Scopus & \emph{citerad}: principen etablerad \\
IMPLEMENTING DIALOGUE INDEPENDENCE & 1995 & WoS, IEEE & \emph{citerad}: kostnader, brister och kritik \\
Mobile GUI Testing Fragility: A Study on Open-Source Android Applications & 2019 & IEEE & \emph{citerad}: att pröva genom gränssnittet \\
Model-view-controller and object teams & 2003 & OA & \emph{citerad}: kostnader, brister och kritik \\
MVP: Model-View-Presenter The Taligent Programming Model for C++ and Java & 1996 & OA & \emph{citerad}: principen etablerad \\
MVVM Revisited: Exploring Design Variants of the Model-View-ViewModel Pattern & 2025 & Scopus & \emph{citerad}: kostnader, brister och kritik \\
N degrees of separation & 1999 & OA & \emph{citerad}: principen etablerad \\
Presenter First: organizing complex GUI applications for test-driven development & 2006 & IEEE, WoS & \emph{citerad}: mönstrens prövbarhet \\
Refactoring model-view-controller & 2005 & OA & \emph{citerad}: principen etablerad \\
Studying the fault-detection effectiveness of GUI test cases for rapidly evolving software & 2005 & IEEE & \emph{citerad}: att pröva genom gränssnittet \\
The abstraction-link-view paradigm & 1992 & OA & \emph{citerad}: kostnader, brister och kritik \\
Usability-supporting architectural patterns & 2004 & IEEE, WoS, Scopus & \emph{citerad}: principen etablerad \\
User interface management systems: themes and variations a review of the literature & 1991 & IEEE & \emph{citerad}: kostnader, brister och kritik \\
Window user interfaces and software maintenance & 1991 & Scopus & \emph{citerad}: principen etablerad \\
{}[Poster] View management for webized mobile AR contents & 2014 & IEEE & stöder påståendet (0.73) \\
A Case for GUI Testing Using Symbolic Execution Poster Abstract & 2007 & IEEE & stöder påståendet (0.82) \\
A case study of user interface management system development and application & 1989 & Scopus, OA & stöder påståendet (0.78) \\
A component- and message-based architectural style for GUI software & 1995 & OA & stöder påståendet (0.95) \\
A Framework for Monkey GUI Testing & 2016 & OA & stöder påståendet (0.95) \\
A Graph Configuration Method Based on UI Service Composition & 2011 & IEEE & stöder påståendet (0.95) \\
A java-based approach for developing Web application system & 1999 & IEEE & stöder påståendet (0.92) \\
A knowledge-based user interface management system & 1988 & OA & stöder påståendet (0.83) \\
A linked data model-view-* approach for decoupled client-server applications & 2020 & Scopus & stöder påståendet (0.72) \\
A logical theory of interfaces and objects & 2002 & IEEE & stöder påståendet (0.98) \\
A model-view separation architecture for GUI application components & 2005 & IEEE & stöder påståendet (0.99) \\
A Model-View-ViewModel (MVVM) Application Framework for Hearing Impairment Diagnosis & 2019 & OA & stöder påståendet (0.99) \\
A Model-View-ViewModel (MVVM) Application Framework for Hearing Impairment Diagnosis - Class Inheritance Architecture & 2020 & IEEE & stöder påståendet (0.88) \\
A Model-View-ViewModel (MVVM) Application Framework for Hearing Impairment Diagnosis - Design and Features & 2020 & IEEE & stöder påståendet (0.90) \\
A Model-View-ViewModel (MVVM) Application Framework for Hearing Impairment Diagnosis - Type Dependency Architecture & 2020 & IEEE & stöder påståendet (0.88) \\
A modern approach to QENS data analysis in Mantid & 2019 & WoS & stöder påståendet (0.98) \\
A Multi-platform Compatible and Open Software Architecture for Computer Numerical Control Systems & 2012 & WoS & stöder påståendet (0.98) \\
A Multidimensional Pivot Table Model Based on MVVM Pattern for Rich Internet Application & 2012 & IEEE & stöder påståendet (0.86) \\
A new architecture model for groupware & 2005 & WoS & stöder påståendet (0.90) \\
A new model for handling input & 1990 & OA, WoS & stöder påståendet (0.98) \\
A Novel Approach for Developing Web-Based Enterprise Information Systems & 2008 & IEEE & stöder påståendet (0.93) \\
A rule-based approach to Web-based application development & 2006 & IEEE, WoS & stöder påståendet (0.96) \\
A Scalable Graphics User Interface Architecture for CNC Application based - on WPF and MVVM & 2011 & WoS & stöder påståendet (0.98) \\
A Simple Software Development Methodology Based on MVP for Android Applications in a Classroom Context & 2015 & IEEE & stöder påståendet (0.75) \\
A Study of Cockpit HMI Simulation Design Based on the Concept of MVC Design Pattern & 2018 & WoS & stöder påståendet (0.85) \\
A Systematic Literature Review on Graphical User Interface Testing Through Software Patterns & 2025 & Scopus & stöder påståendet (0.68) \\
A user interface management system & 1982 & OA & stöder påståendet (0.96) \\
A Web-Based Mashup Environment for On-the-Fly Service Composition & 2008 & IEEE & stöder påståendet (0.85) \\
Abstract interaction tools: a language for user interface management systems & 1988 & OA & stöder påståendet (0.89) \\
Abstraction-link-view paradigm: using constraints to connect user interfaces to applications & 1992 & Scopus & stöder påståendet (0.91) \\
ADVANCED USER INTERFACE GENERATION IN THE SOFTWARE FRAMEWORK FOR MAGNETIC MEASUREMENTS AT CERN & 2010 & WoS & stöder påståendet (0.97) \\
Ajax in Action & 2005 & OA & stöder påståendet (0.82) \\
An adaptable model for teaching mobile app development & 2016 & IEEE & stöder påståendet (0.93) \\
An education-oriented development framework research for business demand alteration & 2008 & IEEE, WoS & stöder påståendet (0.90) \\
An eventflow model of GUIbased applications for testing & 2007 & OA & stöder påståendet (0.84) \\
An Extensible Heuristic-Based Framework for GUI Test Case Maintenance & 2009 & IEEE & stöder påståendet (0.94) \\
An Industrial Case-Study on GUI Testing With RPA & 2022 & IEEE & stöder påståendet (0.97) \\
An On-the-Fly Approach to Web-Based Service Composition & 2008 & IEEE & stöder påståendet (0.85) \\
Application of MVP Architecture in Reengineering of Legacy Financial System & 2009 & IEEE & stöder påståendet (0.99) \\
Approach to Building a Web-based Expert System Interface and Its Application for Software Provisioning in Clouds & 2015 & WoS & stöder påståendet (0.72) \\
ASP.NET and JSP Frameworks in Model View Controller Implementation & 2006 & IEEE & stöder påståendet (0.98) \\
ATOM: Automatic Maintenance of GUI Test Scripts for Evolving Mobile Applications & 2017 & OA & stöder påståendet (0.94) \\
Audiometry: A model-view-viewmodel (MVVM) application framework for hearing impairment diagnosis & 2020 & OA & stöder påståendet (0.86) \\
Automated Model-Based Test Case Generation for Web User Interfaces (WUI) From Interaction Flow Modeling Language (IFML) Models & 2019 & OA & stöder påståendet (0.88) \\
Automated reporting of GUI design violations for mobile apps & 2018 & OA & stöder påståendet (0.94) \\
Automated System Testing Using Visual GUI Testing Tools: A Comparative Study in Industry & 2012 & IEEE & stöder påståendet (0.92) \\
Automated test oracles for GUIs & 2000 & OA & stöder påståendet (0.94) \\
Automated Test Selection for Android Apps Based on APK and Activity Classification & 2020 & OA & stöder påståendet (0.91) \\
Automated test software: Separated and integrated & 2010 & IEEE & stöder påståendet (0.76) \\
Automated user interface testing for web applications and TestComplete & 2012 & OA & stöder påståendet (0.93) \\
Automatically repairing event sequence-based GUI test suites for regression testing & 2008 & OA & stöder påståendet (0.93) \\
AUTOMATICALLY-GENERATED USER INTERFACES FOR MEASUREMENT SOFTWARE FRAMEWORKS: A CASE STUDY ON MAGNETIC PERMABILITY AT CERN & 2009 & WoS & stöder påståendet (0.97) \\
Bridging the Gap: Empowering Use Cases with Task Models & 2010 & WoS & stöder påståendet (0.83) \\
Building Enterprise Applications with Windows Presentation Foundation and the Model View ViewModel Pattern & 2011 & OA & stöder påståendet (0.98) \\
Chat robots with a business mind & 2001 & WoS & stöder påståendet (0.78) \\
Chiron-1: a user interface development system tailored to software environments & 1991 & IEEE & stöder påståendet (0.96) \\
Clean Android Architecture: Take a layered approach to writing clean, testable, and decoupled Android applications & 2022 & IEEE & stöder påståendet (0.94) \\
Combining HCI techniques for better user interfacing & 1996 & IEEE & stöder påståendet (0.97) \\
Comparative Study of Performance and Productivity of MVC and MVVM design patterns & 2018 & OA & stöder påståendet (0.92) \\
Comparison of JavaServer Pages and XSLT: A software engineering perspective & 2004 & Scopus & stöder påståendet (0.72) \\
Completing the job interface design & 1992 & IEEE & stöder påståendet (0.96) \\
Conceptualization and Evaluation of Component-Based Testing Unified with Visual GUI Testing: An Empirical Study & 2015 & IEEE & stöder påståendet (0.91) \\
Connecting Interaction Models and Application Logic for Model-Driven Generation of Web-Based Graphical User Interfaces & 2013 & IEEE & stöder påståendet (0.95) \\
Creating GUI Testing Tools Using Accessibility Technologies & 2009 & IEEE & stöder påståendet (0.86) \\
Crowdsourcing GUI Tests & 2013 & OA & stöder påståendet (0.97) \\
Declarative interface models for user interface construction tools: the Mastermind approach & 1996 & OA & stöder påståendet (0.78) \\
Design abstractions for innovative web applications & 2007 & WoS & stöder påståendet (0.78) \\
Design alternatives for user interface management sytems based on experience with COUSIN & 1985 & OA & stöder påståendet (0.92) \\
Design and implementation of a run-time architecture for a distributed and object-oriented UIMS & 1994 & Scopus & stöder påståendet (0.72) \\
Design and implementation of ATM simulation system based on MVC pattern & 2010 & IEEE & stöder påståendet (0.96) \\
Design and Implementation of Measurement Software for Aviation Bearing Micro-Structures & 2026 & IEEE & stöder påståendet (0.87) \\
Designing an MVC Model for Rapid Web Application Development & 2014 & OA & stöder påståendet (0.91) \\
Detaching control from data models in model-based generation of user interfaces & 2015 & Scopus & stöder påståendet (0.84) \\
Detection of embedded code smells in dynamic web applications & 2012 & IEEE & stöder påståendet (0.86) \\
DEUCE : A Declarative Framework for Extricating User Interface Concerns & 2007 & WoS, Scopus & stöder påståendet (0.98) \\
DEUCE : Separating Concerns in User Interfaces & 2007 & IEEE & stöder påståendet (0.98) \\
Developing applications with a UIMS & 1994 & Scopus & stöder påståendet (0.76) \\
Developing cost-effective model-based techniques for GUI testing & 2006 & OA & stöder påståendet (0.90) \\
Development of ePPMS for research proposals based on integrated spring and Hibernate framework & 2015 & IEEE & stöder påståendet (0.91) \\
DIALOG MANAGEMENT - SUPPORT FOR DIALOG INDEPENDENCE & 1988 & WoS & stöder påståendet (0.88) \\
Dialogue independence based on a structured UIMS interface & 1993 & Scopus & stöder påståendet (0.93) \\
EHBDroid: Beyond GUI testing for Android applications & 2017 & OA & stöder påståendet (0.96) \\
Element-Aware Fine-Tuning of Vision-Language Models for Cost-Efficient GUI Testing in an Industrial Setting & 2025 & WoS & stöder påståendet (0.82) \\
Encapsulating interactive behaviors & 1989 & OA & stöder påståendet (0.94) \\
End-to-End Testing in Web Environments: Addressing Practical Challenges & 2025 & Scopus & stöder påståendet (0.75) \\
Engineering Configuration Graphical User Interfaces from Variability Models & 2017 & OA & stöder påståendet (0.91) \\
Enriching Use Cases with CTTs & 2010 & IEEE, WoS & stöder påståendet (0.82) \\
Estimating Return on Investment for GUI Test Automation Frameworks & 2019 & IEEE & stöder påståendet (0.96) \\
Evaluation of separated concerns in web-based delivery of user interfaces & 2015 & Scopus & stöder påståendet (0.76) \\
Evolution and Fragilities in Scripted GUI Testing of Android applications & 2017 & OA & stöder påståendet (0.82) \\
Generating highly interactive user interfaces & 1989 & OA & stöder påståendet (0.97) \\
Generating user interface code in a model based user interface development environment & 2000 & OA & stöder påståendet (0.83) \\
Goals and objectives for user interface software & 1987 & OA & stöder påståendet (0.78) \\
Graphic user interface modelling and testing automation & 2011 & OA & stöder påståendet (0.95) \\
GUI Testing for Android Applications: A Survey & 2023 & IEEE & stöder påståendet (0.82) \\
GUI testing for mobile applications: objectives, approaches and challenges & 2020 & OA & stöder påståendet (0.87) \\
GUI Testing Made Easy & 2009 & IEEE & stöder påståendet (0.95) \\
GUI-Guided Test Script Repair for Mobile Apps & 2020 & OA & stöder påståendet (0.94) \\
GUICop: Approach and toolset for specificationbased GUI testing & 2017 & OA & stöder påståendet (0.87) \\
GUICOP: Specification-Based GUI Testing & 2012 & IEEE & stöder påståendet (0.82) \\
Guide to Efficient Software Design: An MVC Approach to Concepts, Structures, and Models & 2020 & Scopus & stöder påståendet (0.94) \\
Guided GUI testing of android apps with minimal restart and approximate learning & 2013 & OA & stöder påståendet (0.92) \\
History, state and future of user interface management systems & 1988 & OA & stöder påståendet (0.96) \\
How a Language-based GUI Generator Can Influence the Teaching of Object-Oriented Programming & 2011 & WoS, Scopus & stöder påståendet (0.76) \\
Human-Computer Interaction - An Overview of Software Architecture & 2019 & IEEE & stöder påståendet (0.83) \\
HutWindows: an improved architecture for a user interface management system & 1988 & IEEE & stöder påståendet (0.97) \\
IAIMS architecture & 1997 & WoS, Scopus & stöder påståendet (0.91) \\
Implementation of MVC (Model-View-Controller) Architectural to Academic Management Information System with Android Platform Base & 2012 & OA & stöder påståendet (0.94) \\
Improving accessibility with user-tailored interfaces & 2009 & WoS & stöder påståendet (0.88) \\
Insights from Building a GUI Testing Tool & 2023 & IEEE & stöder påståendet (0.78) \\
Integrating formal specifications in the development and testing of UIs by formal model--view--controller pattern & 2025 & Scopus & stöder påståendet (0.96) \\
Integrating support for temporal media into an architecture for graphical user interfaces & 1997 & OA & stöder påståendet (0.86) \\
Interactive user interface design: the chimera UIMS & 1989 & IEEE & stöder påståendet (0.98) \\
Interface usage measurements in a user interface management system & 1988 & OA & stöder påståendet (0.94) \\
Introduction to the Serpent User Interface Management System & 1988 & OA & stöder påståendet (0.88) \\
Investigating the adoption and maintenance of web GUI testing: Insights from GitHub repositories & 2025 & OA & stöder påståendet (0.91) \\
ITS - A TOOL FOR RAPIDLY DEVELOPING INTERACTIVE APPLICATIONS & 1990 & WoS, OA & stöder påståendet (0.96) \\
JAutomate: A Tool for System- and Acceptance-test Automation & 2013 & IEEE & stöder påståendet (0.82) \\
Java implementations of user-interface frameworks & 1997 & IEEE, WoS, Scopus & stöder påståendet (0.95) \\
Layered protocols for computer-human dialogue. I: Principles & 1988 & Scopus & stöder påståendet (0.84) \\
LLM-Based Labelling of Recorded Automated GUI-Based Test Cases & 2025 & WoS & stöder påståendet (0.91) \\
LLM-Empowered Scriptless Functional Testing & 2025 & WoS & stöder påståendet (0.90) \\
Lowering Barriers to Application Development With Cloud-Native Domain-Specific Functions & 2022 & WoS & stöder påståendet (0.91) \\
Magni - A Framework for Developing Context-Aware Mobile Applications & 2017 & WoS & stöder påståendet (0.99) \\
Maintenance of Android Widget-Based GUI Testing: A Taxonomy of Test Case Modification Causes & 2018 & IEEE & stöder påståendet (0.97) \\
Maintenance of automated test suites in industry: An empirical study on Visual GUI Testing & 2016 & OA & stöder påståendet (0.91) \\
Making GUI Testing Practical: Bridging the Gaps & 2015 & IEEE & stöder påståendet (0.88) \\
Mastering Windows Presentation Foundation: Master the art of building modern desktop applications on Windows & 2017 & IEEE & stöder påståendet (0.91) \\
Measuring Maintainability of Web Applications Using an Extensible MVC Architecture & 2022 & Scopus & stöder påståendet (0.91) \\
Migration model for rich internet applications based on PureMVC framework & 2010 & IEEE & stöder påståendet (0.93) \\
MIKE: the menu interaction kontrol environment & 1986 & OA & stöder påståendet (0.90) \\
Model View Controller - A Comparative .NET Versus Java Analysis & 2018 & WoS, Scopus & stöder påståendet (0.90) \\
Model-based testing of multiple GUI variants using the GUI test generator & 2010 & OA & stöder påståendet (0.98) \\
Model-View-ViewModel (MVVM) & 2022 & OA & stöder påståendet (0.96) \\
Modern JavaScript frameworks: A Survey Study & 2018 & IEEE & stöder påståendet (0.88) \\
Modernizing JavaServer pages by transformation & 2005 & WoS, Scopus & stöder påståendet (0.99) \\
Multimedia dialogue management in agent-based open service environments & 1999 & WoS, Scopus & stöder påståendet (0.96) \\
MVC Architecture Driven Design and Implementation of Java Framework for Developing Desktop Application & 2014 & OA & stöder påståendet (0.98) \\
MVC Architecture: A Comparative Study Between Laravel Framework and Slim Framework in Freelancer Project Monitoring System Web Based & 2019 & OA & stöder påståendet (0.86) \\
MVC: Model--View--Controller & 2023 & OA & stöder påståendet (0.96) \\
MVCLang: A software modeling language for the model-View-Controller design pattern & 2020 & Scopus & stöder påståendet (0.72) \\
MVVM Demonstration Using C\# WPF & 2023 & OA & stöder påståendet (0.95) \\
O/live: Transparent Distribution, Persistence, and Partial Replication for Ubiquitous User Interfaces & 2014 & WoS & stöder påståendet (0.98) \\
On separation of platform-independent particles in user interfaces: Survey on separation of concerns in user interface design & 2015 & Scopus & stöder påståendet (0.72) \\
Pattern-Oriented Software Architecture Volume 1: A System of Patterns & 1996 & OA & stöder påståendet (0.92) \\
PENERAPAN POLA ARSITEKTUR MVVM PADA PERANCANGAN APLIKASI PENGADUAN MASYARAKAT BERBASIS ANDROID & 2023 & OA & stöder påståendet (0.94) \\
Personal programming and the object computer & 2020 & Scopus & stöder påståendet (0.97) \\
PHP Framework for Database Management Based on MVC Pattern & 2011 & OA & stöder påståendet (0.84) \\
PICASSO: A User Interface Management System for Real-Time Applications & 2008 & OA & stöder påståendet (0.98) \\
Practical Escape of Exploration Tarpits for Mini-Game Testing in an Industrial Setting & 2025 & WoS & stöder påståendet (0.92) \\
Pro WPF and Silverlight MVVM : Effective Application Development with Model-View-ViewModel & 2010 & OA & stöder påståendet (0.98) \\
Programming Situational Mobile Web Applications with Cloud-Mobile Convergence: An Internetware-Oriented Approach & 2019 & IEEE & stöder påståendet (0.72) \\
Research and Design on Library Management System Based on Struts and Hibernate Framework & 2009 & IEEE & stöder påståendet (0.89) \\
Research on MVP design pattern modeling based on MDA & 2020 & Scopus & stöder påståendet (0.82) \\
Research on. NET Intelligent Platform Based on Three Layer Structure & 2023 & IEEE & stöder påståendet (0.96) \\
Reverse Engineering of Event Handlers of RAD-Based Applications & 2011 & IEEE & stöder påståendet (0.96) \\
Scripted and scriptless GUI testing for web applications: An industrial case & 2023 & OA & stöder påståendet (0.91) \\
Scripted GUI Testing of Android Apps & 2017 & OA & stöder påståendet (0.93) \\
Scriptless GUI Testing on Mobile Applications & 2022 & IEEE & stöder påståendet (0.96) \\
Semantic Matching based Test Case Reuse for Android Applications & 2025 & WoS & stöder påståendet (0.88) \\
Semantic User Interfaces & 2012 & WoS & stöder påståendet (0.93) \\
Separating Application Functionality from the User Interface in a Distributed Environment & 1996 & Scopus & stöder påståendet (0.97) \\
Separating graphics from application in the design of user interfaces & 1990 & Scopus & stöder påståendet (0.97) \\
Separating semantics from rendering: A scene graph based architecture for graphics applications & 2011 & Scopus & stöder påståendet (0.88) \\
Separation of Concerns and Linguistic Integration in WebDSL & 2010 & IEEE & stöder påståendet (0.78) \\
Separation of user interface and application with the Delft Direct Manipulation Manager (D2M2) & 1993 & Scopus & stöder påståendet (0.99) \\
Separations of concerns in the Chiron-1 user interface development and management system & 1993 & OA, WoS & stöder påståendet (0.99) \\
Simulation Testing of Maritime CyberPhysical Systems: Application of ModelViewViewModel & 2022 & OA & stöder påståendet (0.96) \\
SITAR: GUI Test Script Repair & 2016 & IEEE & stöder påståendet (0.92) \\
Software Support for User Interface Description Language & 2011 & WoS, Scopus & stöder påståendet (0.88) \\
Studying the Characteristics of a \textquotedbl{}Good\textquotedbl{} GUI Test Suite & 2006 & OA & stöder påståendet (0.97) \\
Supporting parametrization of business games for multiple educational settings & 2007 & IEEE & stöder påståendet (0.95) \\
Surface interaction: A paradigm and model for separating application and interface & 1990 & Scopus & stöder påståendet (0.98) \\
Survey on user interface programming & 1992 & OA & stöder påståendet (0.85) \\
Technical Approach in Reducing GUI Development Cycle Time for Adopting the Connectivity Solutions in Automotive Infotainment Systems & 2014 & WoS & stöder påståendet (0.83) \\
Testing user interfaces in applications & 2008 & IEEE & stöder påståendet (0.98) \\
The engineering information system user interface management system (EIS-UIMS) & 1989 & IEEE & stöder påståendet (0.96) \\
The Exploration and Practice of MVVM Pattern on Android Platform & 2016 & OA & stöder påståendet (0.96) \\
THE GADGETS USER INTERFACE MANAGEMENT-SYSTEM & 1991 & WoS & stöder påståendet (0.97) \\
The Origin of Model-View-Presenter Pattern by Means of System Decomposition & 2025 & Scopus & stöder påståendet (0.96) \\
The Research of PHP Development Framework Based on MVC Pattern & 2009 & IEEE, WoS & stöder påståendet (0.96) \\
The Research of Web Application Framework Based on SSH & 2008 & IEEE & stöder påståendet (0.91) \\
The Role of Model-View Controller in Object Oriented Software Development & 2019 & OA & stöder påståendet (0.96) \\
The University of Alberta user interface management system & 1985 & OA & stöder påståendet (0.97) \\
The VIP interface design system & 1988 & IEEE & stöder påståendet (0.96) \\
Tools for building advanced user interfaces & 1986 & IEEE & stöder påståendet (0.97) \\
Toward multi-disciplinary model-based (re)design of sustainable user interfaces & 2008 & WoS & stöder påståendet (0.93) \\
Towards a comprehensive user interface management system & 1983 & OA & stöder påståendet (0.88) \\
TOWARDS A NEW QUALITY OF AUTOMATION IN COMPLEX MAN-MACHINE SYSTEMS & 1992 & WoS, Scopus & stöder påståendet (0.72) \\
Towards High Quality Mobile Applications: Android Passive MVC Architecture & 2014 & OA & stöder påståendet (0.96) \\
Towards Smart User Interface Design & 2012 & IEEE & stöder påståendet (0.88) \\
Towards Unit Testing of User Interface Code for Android Mobile Applications & 2011 & WoS & stöder påståendet (0.97) \\
Unified User Interface for a CAD System & 1985 & IEEE & stöder påståendet (0.76) \\
Unit testing the UI & 2007 & WoS & stöder påståendet (0.99) \\
User Interface & 2016 & WoS & stöder påståendet (0.91) \\
ViewModel visualization tool & 2016 & WoS, Scopus & stöder påståendet (0.97) \\
Views and patterns in e-commerce application design & 2002 & WoS & stöder påståendet (0.94) \\
Visual GUI testing in practice: challenges, problemsand limitations & 2014 & OA & stöder påståendet (0.95) \\
VISUAL PROGRAMMING INTERFACE FOR AN IMAGE-PROCESSING ENVIRONMENT & 1994 & WoS & stöder påståendet (0.97) \\
WEBSITE E-COMMERCE MENGGUNAKAN MODEL VIEW CONTROLLER ( MVC ) DENGAN FRAMEWORK CODEIGNITER Studi Kasus : Toko Miniatur & 2015 & OA & stöder påståendet (0.86) \\
What test oracle should I use for effective GUI testing? & 2003 & IEEE & stöder påståendet (0.95) \\
Who's to Blame? Rethinking the Brittleness of Automated Web GUI Testing from a Pragmatic Perspective & 2025 & WoS & stöder påståendet (0.98) \\
Why many challenges with GUI test automation (will) remain & 2021 & OA & stöder påståendet (0.91) \\
WISBuilder: A Framework for Facilitating Development of Web-Based Information Systems & 2006 & IEEE & stöder påståendet (0.91) \\
XIS-UML Profile for eXtreme Modeling Interactive Systems & 2007 & IEEE & stöder påståendet (0.78) \\
A critical examination of separable user interface management systems & 1997 & OA & kvalificerar eller motsäger påståendet (0.91) \\
A Failed attempt at creating Guidelines for Visual GUI Testing: An industrial case study & 2021 & WoS & kvalificerar eller motsäger påståendet (0.94) \\
A Testing Method for Detecting Data Inconsistencies in Model View Controller Architecture & 2022 & IEEE & kvalificerar eller motsäger påståendet (0.97) \\
A Validated Set of Smells in Model-View-Controller Architectures & 2016 & IEEE & kvalificerar eller motsäger påståendet (0.98) \\
An Adaptive Development Framework for Web-Based Enterprise Information System & 2008 & IEEE & kvalificerar eller motsäger påståendet (0.90) \\
An empirical catalog of code smells for the presentation layer of Android apps & 2019 & OA & kvalificerar eller motsäger påståendet (0.94) \\
An Initial Characterization of Industrial Graphical User Interface Systems & 2009 & IEEE & kvalificerar eller motsäger påståendet (0.91) \\
Architectural Patterns in Android Development: Comparing MVP, MVVM, and MVI & 2024 & OA & kvalificerar eller motsäger påståendet (0.97) \\
Clean Architecture: Impact on Performance and Maintainability of Native Android Projects & 2024 & Scopus & kvalificerar eller motsäger påståendet (0.83) \\
Dynamically choosing the appropriate input control strategy using novel Aspect-Oriented Software Development approach & 2012 & WoS & kvalificerar eller motsäger påståendet (0.87) \\
Interaction ergonomics, control and separation: open problems in user interface management & 1987 & Scopus & kvalificerar eller motsäger påståendet (0.95) \\
Interacto: A Modern User Interaction Processing Model & 2022 & Scopus & kvalificerar eller motsäger påståendet (0.76) \\
Light-Weight and Scalable Hierarchical-MVC Architecture for Cloud Web Applications & 2019 & IEEE & kvalificerar eller motsäger påståendet (0.91) \\
MAMP: A design model for object-oriented visualization systems & 1999 & WoS & kvalificerar eller motsäger påståendet (0.82) \\
MVVM in the Era of Modern Android: A Systematic Literature Review and Taxonomy of Architectural Trade-offs & 2026 & Scopus & kvalificerar eller motsäger påståendet (0.84) \\
Naked objects & 2002 & Scopus & kvalificerar eller motsäger påståendet (0.76) \\
On the long-term use of visual gui testing in industrial practice: a case study & 2017 & OA & kvalificerar eller motsäger påståendet (0.94) \\
Rethinking MVC Violations as Security-Aware Architectural Smells in Android Applications & 2026 & Scopus & kvalificerar eller motsäger påståendet (0.91) \\
Review of iOS architectural pattern for testability, modifiability, and performance quality & 2019 & OA & kvalificerar eller motsäger påståendet (0.99) \\
Revisiting the notion of GUI testing & 2019 & Scopus & kvalificerar eller motsäger påståendet (0.63) \\
Taligent MVP in interactive statistical graphics & 2008 & WoS & kvalificerar eller motsäger påståendet (0.97) \\
TOOLS FOR BUILDING THE HUMAN-COMPUTER INTERFACE OF A DECISION-SUPPORT SYSTEM & 1995 & WoS, Scopus & kvalificerar eller motsäger påståendet (0.98) \\
TOWARDS A COMMON CONTEXT FOR USER-INTERFACE MANAGEMENT-SYSTEM DESIGN & 1992 & WoS & kvalificerar eller motsäger påståendet (0.94) \\
Using active data in a UIMS & 1988 & OA & kvalificerar eller motsäger påståendet (0.84) \\
Visual GUI Testing: Automating High-level Software Testing in Industrial Practice & 2015 & OA & kvalificerar eller motsäger påståendet (0.96) \\
\end{longtable}
\endgroup

